\subsection{\texorpdfstring{映射 \(\Theta_{a}\) 的性质}{映射 \textbackslash Theta\_\{a\} 的性质}}

回顾（VAR, R2, p.~46 与 p.~47, 11.1.3, d)）：若 K 是 C 的一个紧子集，则存在由 K 的紧邻域组成的基本系，其中的邻域都是 C 的块（也就是说，对每个 \(x \in V\)，存在 C 在 x 处的一个卡 \((\mathrm{U}, \varphi, \mathbf{C})\)，使得 \(\varphi(\mathrm{U} \cap \mathrm{K})\) 是 C 的某个闭半空间的开子集）。我们用 \(\partial V\) 表示赋予由 C 的定向所诱导的定向后的块 V 的边界 (VAR, R2, p.~47)，并用 dz 表示内射 \(\partial V \to C\) 的微分。

命题 1. --- 设 a 是 A 的一个元素，U 是 Sp(a) 的一个开邻域，并设 \(f \in \mathcal{O}(\mathrm{U};\mathrm{A})\)。设 V 是 Sp(a) 的一个包含于 U 的紧邻域，且 V 是 C 的一个块。

则映射 \(z \mapsto f(z)(z - a)^{-1}\) 在 \(\partial\mathrm{V}\) 上连续；微分形式 \(f(z)(z - a)^{-1}dz\) 在 \(\partial\mathrm{V}\) 上可积，并且有

\[
\Theta_ {a} (f) = \frac {1}{2 i \pi} \int_ {\partial \mathrm{V}} f (z) (z - a) ^ {- 1} d z.
\]

设 \(h\) 是从 \(\mathbf{C}\) 到 \(\mathbf{C}\) 中的无穷次可微映射，它在 \(\mathrm{Sp}(a)\) 的一个邻域中等于 1，并且其紧支集包含于 \(V\) 的内部 (I, p.~52, 引理 1)。设 \(u\) 是从 \(\mathbf{C}\) 到 A 中的映射，使得二元组 \((h, u)\) 适配于 \(a\)（参见 I, p.~53, 例 3）。相伴的微分形式是 \(\omega = du \wedge dz\)。由此得 \(f\omega = fdu \wedge dz = d(fu dz)\)，因为 \(f\) 是全纯的。此外，\(u(z) = (z - a)^{-1}\) 在 \(V\) 的边界上成立。再者，微分形式 \(fu dz\) 在 U 上是 \(C^1\) 类的。于是，

\[
2 i \pi \Theta_ {a} (f) = \int_ {\mathrm{V}} d (f u d z) = \int_ {\partial \mathrm{V}} f u d z = \int_ {\partial \mathrm{V}} f (z) (z - a) ^ {- 1} d z
\]

这是由公式 (5) 以及关于块 V 的 Stokes 公式 (VAR, R2, p.~47, 11.2.3) 得到的。

\[
\text { COROLLARY. } - \text {   Let   } a \in \mathrm{A}. \text {   One has   } \Theta_ {a} (1) = 1.
\]

设 \(\mathrm{R} > \varrho(a)\) 是一个实数。设 V 是以 0 为中心、以 R 为半径的闭圆盘，于是 \(\mathrm{Sp}(a) \subset \mathring{\mathrm{V}}\)。它是 C 的一个块，其边界 \(\partial\mathrm{V}\) 是以 0 为中心、以 R 为半径的圆周。对 \(z \in \mathbf{C} - \mathring{\mathrm{V}}\)，我们有公式 \((z - a)^{-1} = z^{-1}(1 - z^{-1}a)^{-1} = \sum_{j=1}^{+\infty} z^{-j} a^{j-1}\)。该级数对 \(z \in \partial V\) 一致收敛。因此我们有

\[
\Theta_ {a} (1) = \frac {1}{2 i \pi} \sum_ {j = 1} ^ {+ \infty} a ^ {j - 1} \int_ {\partial \mathrm{V}} z ^ {- j} d z = 1
\]

因为

\[
\int_ {\partial \mathrm{V}} z ^ {j} d z = 0
\]

对每个整数 \(j \neq -1\) 成立，且

\[
\int_ {\partial \mathrm{V}} z ^ {- 1} d z = 2 i \pi
\]

(VAR, R2, p.~44, 10.4.5 及 p.~47, 11.2.1, 例)。

引理 8. --- 设 U 是 \(C^{n}\) 的一个开子集。设 \(\omega_{1}\) 是 U 上的一个 n 次连续微分形式，具有紧支集且取值于 A（相应地，设 \(\omega_{2}\) 是 C 上的一个 2 次连续微分形式，具有紧支集且取值于 C）。我们用 \(\pi_{1}\) 与 \(\pi_{2}\) 表示 U × C 到 U 与 C 上的典范投影。U × C 上的微分形式 \(\pi_{1}^{*}\omega_{1} \wedge \pi_{2}^{*}\omega_{2}\) 连续、具有紧支集且取值于 A。我们有

\[
\int_ {\mathrm{U} \times \mathbf {C}} \pi_ {1} ^ {*} \omega_ {1} \wedge \pi_ {2} ^ {*} \omega_ {2} = \left(\int_ {\mathbf {C}} \omega_ {2}\right) \left(\int_ {\mathrm{U}} \omega_ {1}\right).
\]

设 \(\mu_{n}\) 是 \(C^{n}\) 上的 Lebesgue 测度，\(\mu_{1}\) 是 C 上的 Lebesgue 测度。存在 \(\psi_{1} \in \mathcal{K}(U;A)\) 与 \(\psi_{2} \in \mathcal{K}(C)\)，使得与 \(\omega_{1}\) 相伴的向量测度等于 \(\psi_{1} \cdot \mu_{n}\)，而与 \(\omega_{2}\) 相伴的向量测度等于 \(\psi_{2} \cdot \mu_{1}\)。与微分形式 \(\pi_{1}^{*}\omega_{1} \wedge \pi_{2}^{*}\omega_{2}\) 相伴的向量测度是 \((\psi_{1} \otimes \psi_{2}) \cdot \mu_{n} \otimes \mu_{1}\)。

设 \(\ell\) 是 A 上的一个连续线性形式。由 INT, VI, §2, 第 2 小节, 定义 2 以及乘积测度的定义 (INT, III, §4, 第 1 小节, 定义 1) 得

\[
\begin{array}{r l} & {\ell \Big (\int_ {\mathrm{U} \times \mathbf {C}} \pi_ {1} ^ {*} \omega_ {1} \wedge \pi_ {2} ^ {*} \omega_ {2} \Big) = \int_ {\mathrm{U} \times \mathbf {C}} \ell \circ (\psi_ {1} \otimes \psi_ {2}) \mu_ {n} \otimes \mu_ {1}} \\ & {\qquad = \int_ {\mathrm{U} \times \mathbf {C}} \psi_ {2} (z) \ell (\psi_ {1} (x)) d \mu_ {n} (x) d \mu_ {1} (z)} \\ & {\qquad = \Big (\int_ {\mathbf {C}} \psi_ {2} (z) d \mu_ {1} (z) \Big) \Big (\int_ {\mathrm{U}} \ell (\psi_ {1} (x)) d \mu_ {n} (x) \Big)} \\ & {\qquad = \Big (\int_ {\mathbf {C}} \psi_ {2} \mu_ {1} \Big) \ell \Big (\int_ {\mathrm{U}} \psi_ {1} \mu_ {n} \Big),} \end{array}
\]

由此即得结论 (INT VI, 同前)。

引理 9. --- 设 \(\boldsymbol{a} = (a_1, \ldots, a_n) \in \mathbb{A}^n\)。设 \(p \in \mathbf{N}\)，\(\boldsymbol{a}' = (a_{n+1}, \ldots, a_{n+p}) \in \mathbb{A}^p\)。则我们有 \(\Theta_{(\boldsymbol{a}, \boldsymbol{a}')} \circ \pi_{n,n+p}^* = \Theta_{\boldsymbol{a}}\)。特别地，我们有 \(\Theta_{\boldsymbol{a}}(1) = 1\)。

由于 \(\pi_{n,n+p} = \pi_{n,n+1} \circ \cdots \circ \pi_{n+p-1,n+p}\)，只需对 \(p = 1\) 的情形证明第一个断言，我们现在就作此假设。我们简记 \(\pi = \pi_{n,n+1}\)。于是只需证明：对 \(\mathrm{Sp}^n(\boldsymbol{a})\) 的每个开邻域 U 以及每个函数 \(f \in \mathcal{O}(\mathrm{U};\mathrm{A})\)，有 \(\Theta_{(\boldsymbol{a},a_{n+1})}(f \circ \pi) = \Theta_{\boldsymbol{a}}(f)\)。记 \(g = f \circ \pi\)。设 \(h\)（相应地，\(h'\)）是从 \(\mathbf{C}^n\) 到 \(\mathbf{C}\) 中（相应地，从 \(\mathbf{C}\) 到 \(\mathbf{C}\) 中）的无穷次可微映射，它在 \(\mathrm{Sp}^n(\boldsymbol{a})\)（相应地，\(\mathrm{Sp}(\boldsymbol{a}')\)）的一个邻域中等于 1，且具有包含于 U（相应地，\(\mathbf{C}\)）中的紧支集。存在映射 ( \(u_1, \ldots, u_n\) )：从 \(\mathbf{C}^n\) 到 A 中、无穷次可微，使得序列 ( \(h, u_1, \ldots, u_n\) ) 适配于 \(\boldsymbol{a}\) (I, p.~54, 引理 3)；又存在无穷次可微映射 \(u_{n+1}\)：从 \(\mathbf{C}\) 到 A 中，使得二元组 ( \(h', u_{n+1} \) ) 适配于 \(a_{n+1}\)（同前）。

对 \(z \in C^{n}\) 与 \(z_{n+1} \in C\)，令 \(h''(z, z_{n+1}) = h(z)h'(z_{n+1})\)，\(u_{n+1}''(z, z_{n+1}) = h(z)u_{n+1}(z_{n+1})\)。函数 \(h''\) 与 \(u_{n+1}''\) 在 \(C^{n+1}\) 中无穷次可微。函数 \(h''\) 在 \(\mathrm{Sp}^{n+1}(a, a_{n+1})\) 的一个邻域中等于 1，且具有包含于 \(U \times C\) 的紧支集。对每个 \(\boldsymbol{w} = (z, z_{n+1}) \in \mathbf{C}^{n+1}\)，有

\[
\begin{array}{c} (z _ {1} - a _ {1}) (u _ {1} \circ \pi) (\boldsymbol {w}) + \dots + (z _ {n} - a _ {n}) (u _ {n} \circ \pi) (\boldsymbol {w}) \\ + (z _ {n + 1} - a _ {n + 1}) u _ {n + 1} ^ {\prime \prime} (\boldsymbol {w}) = 1 - h (\boldsymbol {z}) + h (\boldsymbol {z}) (1 - h ^ {\prime} (z _ {n + 1})) \\ = 1 - h ^ {\prime \prime} (\boldsymbol {w}), \end{array}
\]

这表明序列 \((h''_{1}, u_{1} \circ \pi, \ldots, u_{n} \circ \pi, u_{n+1}')\) 适配于 \((\boldsymbol{a}, a_{n+1})\)。设 \(\omega\) 是相伴的微分形式。

微分形式 \(du_{1} \wedge dz_{1} \wedge \cdots \wedge du_{n} \wedge dz_{n} \wedge dh\)（它定义在 \(C^{n}\) 上）的次数为 \(2n + 1\)，故为零。因此，

\[
\begin{array}{l} \omega = d (u _ {1} \circ \pi) \wedge d z _ {1} \wedge \dots \wedge d (u _ {n} \circ \pi) \wedge d z _ {n} \wedge d u _ {n + 1} ^ {\prime \prime} \wedge d z _ {n + 1} \\ = (h \circ \pi) d (u _ {1} \circ \pi) \wedge d z _ {1} \wedge \dots \wedge d (u _ {n} \circ \pi) \wedge d z _ {n} \wedge d u _ {n + 1} \wedge d z _ {n + 1}. \end{array}
\]

由于 \(g = f \circ \pi\)，公式 (5) 与引理 8 蕴涵

\[
\begin{array}{c} \Theta_ {\boldsymbol {a}, \boldsymbol {a} ^ {\prime}} (g) = \frac {(n + 1) !}{(2 i \pi) ^ {n + 1}} \int_ {\mathrm{U} \times \mathbf {C}} g \omega = \frac {(n + 1) !}{(2 i \pi) ^ {n + 1}} \Big (\int_ {\mathbf {C}} d u _ {n + 1} \wedge d z _ {n + 1} \Big) \\ \qquad \qquad \times \Big (\int_ {\mathrm{U}} f h   d u _ {1} \wedge d z _ {1} \wedge \dots \wedge d u _ {n} \wedge d z _ {n} \Big). \end{array}
\]

一方面，我们有

\[
\int_ {\mathbf {C}} d u _ {n + 1} \wedge d z _ {n + 1} = 2 i \pi \Theta_ {a _ {n + 1}} ^ {\mathbf {C}} (1) = 2 i \pi \cdot 1
\]

（由推论 5）。另一方面，I, p.~54, 引理 4 的 c) 以及闭形式的积分为零这一事实 (VAR, R2, p.~48, 11.2.4) 蕴涵

\[
\begin{array}{c} (n + 1) \int_ {\mathrm{U}} f h d u _ {1} \wedge d z _ {1} \wedge \dots \wedge d u _ {n} \wedge d z _ {n} = \\ \int_ {\mathrm{U}} f   d u _ {1} \wedge d z _ {1} \wedge \dots \wedge d u _ {n} \wedge d z _ {n} = \frac {(2 i \pi) ^ {n}}{n !} \Theta_ {\boldsymbol {a}} (f). \end{array}
\]

这样就得到

\[
\Theta_ {\boldsymbol {a}, \boldsymbol {a} ^ {\prime}} (g) = \frac {(n + 1) !}{(2 i \pi) ^ {n + 1}} \times \frac {(2 i \pi) ^ {n}}{(n + 1) !} \Theta_ {\boldsymbol {a}} (f) \times 2 i \pi = \Theta_ {\boldsymbol {a}} (f).
\]

最后，公式 \(\Theta_{\boldsymbol{a}}(1)=1\) 由前述结果以及命题 1 的推论得出。

引理 10. --- 设 \(\boldsymbol{a} \in \mathrm{A}^n\)。设 \(g\) 是 \(\mathbf{C}^n\) 上的一个以 A 为系数的多项式函数，并设 \(f \in \mathcal{O}(\mathrm{Sp}^n(\boldsymbol{a});\mathrm{A})\)。我们有 \(\Theta_{\boldsymbol{a}}(gf) = g(\boldsymbol{a})\Theta_{\boldsymbol{a}}(f)\)。特别地，我们有 \(\Theta_{\boldsymbol{a}}(g) = g(\boldsymbol{a})\)。

由引理 9，只需证明第一个断言。

我们用 \(z_{1},\ldots,z_{n}\) 表示 \(C^{n}\) 上的坐标函数。由于映射 \(\Theta_{a}\) 是 A-线性的，只需对 \(g=z_{1}^{e_{1}}\cdots z_{n}^{e_{n}}\)（其中 \((e_{1},\ldots,e_{n})\in\mathbf{N}^{n}\)）的情形证明引理的断言。对 \(e_{1}+\cdots+e_{n}\) 作归纳，可化归为存在整数 i 使得 \(1\leqslant i\leqslant n\) 且 \(g=z_{i}\) 的情形。

设 U 是 \(\mathrm{Sp}^{n}(\boldsymbol{a})\) 的一个开邻域。设 \((h,u_{1},\ldots,u_{n})\) 是适配于 \(\boldsymbol{a}\) 的序列，且 h 的支集包含于 U (I, p.~52, 引理 1 及 I, p.~54, 引理 3)，并设 \(\omega\) 是相伴的微分形式。由 I, p.~54, 引理 4, a)，存在一个微分形式 \(\beta\)，使得

\[
(z _ {i} - a _ {i}) \omega = d (h \beta \wedge d z _ {1} \wedge \dots \wedge d z _ {n}).
\]

因此，对每个函数 \(f \in \mathcal{O}(\mathrm{U};\mathrm{A})\)，有

\[
\left(z _ {i} - a _ {i}\right) f \omega = f d \left(h \beta \wedge d z _ {1} \wedge \dots \wedge d z _ {n}\right) = d \left(f h \beta \wedge d z _ {1} \wedge \dots \wedge d z _ {n}\right)
\]

这是因为 f 全纯，从而 \(df \wedge dz_{1} \wedge \cdots \wedge dz_{n} = 0\)。应用 Stokes 公式 (VAR, R2, p.~48, 11.2.4)，我们得到 \(\int_{\mathrm{U}}(z_{i} - a_{i})f\omega = 0\)，由此得

\[
\Theta_ {\boldsymbol {a}} (z _ {i} f) = \frac {n !}{(2 i \pi) ^ {n}} \int_ {\mathrm{U}} z _ {i} f \omega = \frac {n !}{(2 i \pi) ^ {n}} \int_ {\mathrm{U}} a _ {i} f \omega = a _ {i} \Theta_ {\boldsymbol {a}} (f)
\]

（由公式 (5)）。结论得证。

命题 2. — 设 \(\varrho_{1},\ldots,\varrho_{n}\) 是 \(>0\) 的实数，并设 \(\mathrm{U}\subset\mathbf{C}^{n}\) 是以 0 为中心、以 \(\varrho_{i}\) 为半径的开圆盘的乘积多圆柱。设

\[
\sum_ {(k _ {1}, \dots , k _ {n}) \in \mathbf {N} ^ {n}} c (k _ {1}, \dots , k _ {n}) \mathrm{X} _ {1} ^ {k _ {1}} \dots \mathrm{X} _ {n} ^ {k _ {n}} \in \mathrm{A} [ [ \mathrm{X} _ {1}, \dots , \mathrm{X} _ {n} ] ]
\]

是一个以 A 为系数的形式级数。假设该级数在 U 中收敛，并用 f 表示作为其和的 U 中的全纯函数。

设 \(\boldsymbol{a}=(a_{1},\ldots,a_{n})\in\mathrm{A}^{n}\) 满足 \(\varrho(a_{i})<\varrho_{i}\)（对 \(1\leqslant i\leqslant n\)）。则 \(\mathrm{Sp}^{n}(\boldsymbol{a})\subset\mathrm{U}\)，A 的元素所成的族 \((c(k_{1},\ldots,k_{n})a_{1}^{k_{1}}\cdots a_{n}^{k_{n}})\) 绝对可和，且

\[
\Theta_ {\boldsymbol {a}} (f) = \sum_ {(k _ {1}, \dots , k _ {n}) \in \mathbf {N} ^ {n}} c (k _ {1}, \dots , k _ {n}) a _ {1} ^ {k _ {1}} \dots a _ {n} ^ {k _ {n}}.
\]

对 A 的每个特征标 \(\chi\) 以及每个整数 \(i\)（满足 \(1 \leqslant i \leqslant n\)），有 \(|\chi(a_i)| \leqslant \varrho(a_i) < \varrho_i\)，于是由联合谱的定义得 \(\mathrm{Sp}^n(\boldsymbol{a}) \subset \mathrm{U}\)。设 \(z_1, \ldots, z_n\) 是 \(\mathbf{C}^n\) 的坐标函数在 U 上的限制。则族 \((c(k_1, \ldots, k_n)z_1^{k_1} \cdots z_n^{k_n})\) 在 \(\mathcal{O}(\mathrm{U}; \mathrm{A})\) 中可和，且和为 \(f\)。顾及引理 10 以及映射 \(\Theta_{\boldsymbol{a}}^{\mathrm{U}}\) 的连续性，族 \((c(k_1, \ldots, k_n)a_1^{k_1} \ldots a_n^{k_n})\) 因而在 A 中可和，且和为 \(\Theta_{\boldsymbol{a}}(f)\)。对 \(1 \leqslant i \leqslant n\)，设 \(\lambda_i\) 是满足 \(\varrho(a_i) < \lambda_i < \varrho_i\) 的实数。存在 \(\mathrm{M}_i < +\infty\)，使得 \(\|a_i^k\| \leqslant \mathrm{M}_i\lambda_i^k\)（对每个整数 \(k \geqslant 0\)）。于是我们有

\[
\sum_ {(k _ {1}, \dots , k _ {n}) \in \mathbf {N} ^ {n}} \| c (k _ {1}, \dots , k _ {n}) \| \| a _ {1} ^ {k _ {1}} \dots a _ {n} ^ {k _ {n}} \| \leqslant
\]

\[
\mathrm{M} _ {1} \dots \mathrm{M} _ {n} \sum_ {(k _ {1}, \dots , k _ {n}) \in \mathbf {N} ^ {n}} \| c (k _ {1}, \dots , k _ {n}) \| \lambda_ {1} ^ {k _ {1}} \dots \lambda_ {n} ^ {k _ {n}}
\]

它依假设是有限的，故族 \((c(k_{1},\ldots,k_{n})a_{1}^{k_{1}}\cdots a_{n}^{k_{n}})\) 绝对可和。

推论. --- 假设 A 非零。设 \(a \in C^{n} \subset A^{n}\)。我们有 \(\operatorname{Sp}_{\mathrm{A}}^{n}(\boldsymbol{a}) = \{\boldsymbol{a}\}\)。对每个芽 \(f \in \mathcal{O}(\{\boldsymbol{a}\}; \mathrm{A})\)，有 \(\Theta_{\boldsymbol{a}}(f) = f(\boldsymbol{a})\)。

命题 3. --- 设 B 是一个交换幺 Banach 代数，\(\varphi\) 是从 A 到 B 中的连续保单位元态射。设 \(\boldsymbol{a} = (a_1, \ldots, a_n) \in \mathrm{A}^n\)。令 \(\boldsymbol{b} = (\varphi(a_1), \ldots, \varphi(a_n))\)，于是 \(\mathrm{Sp}_{\mathrm{B}}^{n}(\boldsymbol{b}) \subset \mathrm{Sp}_{\mathrm{A}}^{n}(\boldsymbol{a})\)。对每个 \(f \in \mathcal{O}(\mathrm{Sp}_{\mathrm{A}}^{n}(\boldsymbol{a}); \mathrm{A})\)，有

\[
\varphi (\Theta_ {\boldsymbol {a}} (f)) = \Theta_ {\boldsymbol {b}} (\varphi_ {*} (f)),
\]

其中 \(\varphi_{*}(f)\) 表示 \(\varphi \circ f\) 在 \(\mathrm{Sp}_{\mathrm{B}}^{n}(\boldsymbol{b})\) 的一个邻域中的芽。

只需证明：对 \(\mathrm{Sp}_{\mathrm{A}}^{n}(\boldsymbol{a})\) 的每个开邻域 U 以及每个 \(f \in \mathcal{O}(\mathrm{U}; \mathrm{A})\)，有 \(\varphi(\Theta_{\boldsymbol{a}}(f)) = \Theta_{\boldsymbol{b}}(\varphi \circ f)\)，其中 \(\varphi \circ f \in \mathcal{O}(\mathrm{U}; \mathrm{B})\)。设 \((h, u_{1}, \ldots, u_{n})\) 是适配于 \(\boldsymbol{a}\) 的序列，其中 h 的支集包含于 U (I, p.~52, 引理 1 及 I, p.~54, 引理 3)。用 \(\omega\) 表示相伴的微分形式。对每个 \(z \in C^{n}\)，有

\[
\sum_ {j = 1} ^ {n} (z _ {j} - b _ {j}) \varphi (u _ {i} (\boldsymbol {z})) = \varphi \left(\sum_ {j = 1} ^ {n} (z _ {j} - a _ {j}) u _ {j} (\boldsymbol {z})\right) = 1 - h (\boldsymbol {z}),
\]

因而序列 \((h,\varphi\circ u_{1},\ldots,\varphi\circ u_{n})\) 适配于 b。设 \(\omega'\) 是相伴的微分形式。用 \(\mu\) 表示 \(C^{n}\) 上的 Lebesgue 测度。设 \(f\in\mathcal{O}(U;A)\)。用 \(\psi\cdot\mu\) 表示与微分形式 \(f\omega\) 相伴的向量测度。与微分形式

\[
(\varphi \circ f) \omega^ {\prime} = (\varphi \circ f) d (\varphi \circ u _ {1}) \wedge d z _ {1} \wedge \dots \wedge d (\varphi \circ u _ {n}) \wedge d z _ {n}
\]

相伴的向量测度等于 \((\varphi \circ \psi) \cdot \mu\)。因此，由公式 (5) 以及 INT, VI, §2, 第 2 小节, 命题 2，我们有

\[
\Theta_ {\boldsymbol {b}} (\varphi \circ f) = \frac {n !}{(2 i \pi) ^ {n}} \int_ {\mathrm{U}} (\varphi \circ f) \mu = \frac {n !}{(2 i \pi) ^ {n}} \varphi \left(\int_ {\mathrm{U}} \psi \mu\right) = \varphi (\Theta_ {\boldsymbol {a}} (f)),
\]

这正是所要证明的。

推论 1. --- 设 \(\chi \in \mathrm{X}(\mathrm{A})\)，\(\boldsymbol{a} \in \mathrm{A}^n\)。对每个芽 \(f \in \mathcal{O}(\mathrm{Sp}^n(\boldsymbol{a}))\)，有 \(\chi(\Theta_{\boldsymbol{a}}(f)) = f(\chi(a_1), \ldots, \chi(a_n))\)。

这可由命题 3（应用于连续保单位元态射 \(\chi\colon\mathbf{A}\to\mathbf{C}\) (I, p.~29, 定理 1)）以及命题 2 的推论（应用于 Banach 代数 C）得到。

注记. --- 假设代数 A 是半单的。设 \(\boldsymbol{a} = (a_{1}, \ldots, a_{n}) \in \mathrm{A}^{n}\)。由 I, p.~38, 命题 8，映射 \(\Theta_{\boldsymbol{a}}^{\mathrm{U}}\) 是唯一的映射 \(\varphi\)：从 \(\mathcal{O}(\mathrm{U})\) 到 A 中，使得 \(\chi(\varphi(f)) = f(\chi(a_{1}), \ldots, \chi(a_{n}))\)（对一切 \(\chi \in \mathrm{X}(\mathrm{A})\) 以及每个函数 \(f \in \mathcal{O}(\mathrm{U})\)）。

推论 2. --- 设 p 是一个 ≥ 1 的整数。对每个族 \((f_{1},\ldots,f_{p})\)：由 \(\mathcal{O}(\mathrm{Sp}^{n}(\boldsymbol{a}))\) 的元素组成，有

\[
\operatorname{Sp} ^ {p} \left(\left(\Theta_ {\boldsymbol {a}} \left(f _ {1}\right), \dots , \Theta_ {\boldsymbol {a}} \left(f _ {p}\right)\right)\right) = \left(f _ {1}, \dots , f _ {p}\right) \left(\operatorname{Sp} ^ {n} (\boldsymbol {a})\right).
\]

特别地，对每个 \(f \in \mathcal{O}(\mathrm{Sp}^n(\boldsymbol{a}))\)，有 \(\mathrm{Sp}(\Theta_{\boldsymbol{a}}(f)) = f(\mathrm{Sp}^n(\boldsymbol{a}))\)。

这由推论 1 及联合谱的定义得出。

例. --- 设 A 是由单位圆上具有绝对收敛 Fourier 级数的函数所组成的复 Banach 代数 (I, p.~19, 例 8)。设 \(\varphi \in\) A。设 \(f\) 是 \(\varphi\) 的值集的某个邻域中的全纯函数芽。则 \(\psi = \Theta_{\varphi}(f)\) 是一个绝对收敛 Fourier 级数，它对每个 \(u \in \mathbf{U}\) 满足 \(\psi(u) = f(\varphi(u))\)（推论 1 应用于特征标 \(\varphi \mapsto \varphi(u)\)）。换句话说，单位圆上的函数 \(f \circ \varphi\) 也具有绝对收敛的 Fourier 级数（“P. Lévy 定理”）。这一结果推广了 Wiener 定理 (I, p.~38, 例 4)，后者涉及函数 \(f(z) = 1/z\)（在 \(\mathbf{C} - \{0\}\) 上）当 \(\varphi\) 不取零值时的情形。

\subsection{逼近定理}

在本小节中，A 是一个复交换幺 Banach 代数。

命题 4. --- 设 L 是 \(\mathbf{C}^{n}\) 的一个多项式凸紧子集，U 是 L 的一个开邻域。对每个函数 \(f\in\mathcal{O}(U;A)\)，存在一个由 \(\mathbf{C}^{n}\) 上以 A 为系数的多项式函数组成的序列，它收敛于 \(f|L\)（依 \(\mathcal{C}(L;A)\) 中的范数）。

可以假设 L 非空且 A 非零。设 P（相应地，\(P_{0}\)）是由 \(C^{n}\) 上以 A 为系数（相应地，以 C 为系数）的多项式函数在 L 上的限制所组成的集合。设 B（相应地，\(B_{0}\)）是 P（相应地，\(P_{0}\)）在

\(\mathcal{C}(\mathrm{L};\mathrm{A})\) 中的 Banach 代数闭包。用 \(\iota\) 表示 A 到 B 中由常值函数所组成的赋范子代数上的单射。

设 \(z_{1},\ldots,z_{n}\) 是 \(\mathbf{C}^{n}\) 的坐标函数在 L 上的限制；它们都是 \(\mathrm{B}_{0}\) 的元素，并且令 \(\boldsymbol{z}=(z_{1},\ldots,z_{n})\)，则依 I, p.~47, 命题 15 有 \(\mathrm{Sp}_{\mathrm{B}_{0}}^{n}(\boldsymbol{z})=\mathrm{L}\)。

设 \(f \in \mathcal{O}(\mathrm{U};\mathrm{A})\)。通过与 \(\iota\) 复合，函数 \(f\) 定义了一个元素 \(f_{\mathrm{B}} = \iota \circ f\)（属于 \(\mathcal{O}(\mathrm{U};\mathrm{B})\)）。由于 \(\mathrm{Sp}_{\mathrm{B}}^{n}(\boldsymbol{z}) \subset \mathrm{Sp}_{\mathrm{B}_{0}}^{n}(\boldsymbol{z}) \subset \mathrm{U}\)，可以构造 B 的元素 \(b = \Theta_{\boldsymbol{z}}(f_{\mathrm{B}})\)。设 \(\boldsymbol{w} = (w_{1},\ldots ,w_{n}) \in \mathrm{L}\)，并设 \(\varphi\) 是从 B 到 A 中的连续保单位元态射 \(g \mapsto g(\boldsymbol{w})\)。我们有 \(\varphi \circ \iota = \mathrm{Id}_{\mathrm{A}}\)，因而 \(\varphi \circ f_{\mathrm{B}} = f\)。由于 \(\varphi (z_i) = w_i\)，I, p.~66, 命题 3 蕴涵 \(\varphi (\Theta_z(f_{\mathrm{B}})) = \Theta_w(\varphi \circ f_{\mathrm{B}})\)。因此我们有

\[
b (\pmb {w}) = \varphi (b) = \varphi (\Theta_ {z} (f _ {\mathrm{B}})) = \Theta_ {\pmb {w}} (\varphi \circ f _ {\mathrm{B}}) = \Theta_ {\pmb {w}} (f) = f (\pmb {w}),
\]

（由 I, p.~65, 命题 2 的推论）。于是，我们有 \(f|_{\mathrm{L}} = b\)；特别地，\(f|_{\mathrm{L}}\) 属于 B。这就证明了该命题。

定理 2 (Oka--Weil). --- 设 K 是 \(C^{n}\) 的一个多项式凸紧子集，P 是由 \(C^{n}\) 上以 A 为系数的多项式函数在 K 的某个邻域中的芽所组成的集合。则 P 在 \(\mathcal{O}(K;A)\) 中稠密。更确切地说，\(\mathcal{O}(K;A)\) 的每个元素都是 P 中元素所成序列的极限。

考察 \(\mathcal{O}(\mathrm{K};\mathrm{A})\) 的一个元素，即某个函数 \(f\in \mathcal{O}(\mathrm{U};\mathrm{A})\) 的芽，其中 U 是 K 的一个开邻域。由 I, p.~48, 引理 7，存在 K 的一个包含于 U 且为多项式凸的紧邻域 L。设 V 是 L 的内部；它是 K 的一个邻域。

依前面的命题，存在一个序列 \((\mathrm{P}_{k})\)：它由 \(C^{n}\) 上以 A 为系数的多项式函数组成，且收敛于 \(f|L\)（依 \(\mathcal{C}(L;A)\) 中的范数）。特别地，序列 \((\mathrm{P}_{k})\) 收敛于 \(f|V\)（在 \(\mathcal{O}(V;A)\) 中）。

按 \(\mathcal{O}(\mathrm{K};\mathrm{A})\) 上拓扑的定义（参见 EVT, II, p. 29, 命题 5），从 \(\mathcal{O}(\mathrm{V};\mathrm{A})\) 到 \(\mathcal{O}(\mathrm{K};\mathrm{A})\) 中的典范映射是连续的。因此，函数 \(\mathrm{P}_k\) 在 K 的某个邻域中的芽所成的序列，收敛于 \(f\) 在 K 的邻域中的芽（在空间 \(\mathcal{O}(\mathrm{K};\mathrm{A})\) 中），这就是所要证明的。

推论 1. --- 设 U 是 K 的一个开邻域。设 \(u_{1}\) 与 \(u_{2}\) 是从 \(\mathcal{O}(\mathrm{U};\mathrm{A})\) 到拓扑空间 X 中的、可经由 \(\mathcal{O}(\mathrm{K};\mathrm{A})\) 分解的连续映射。则 \(u_{1}=u_{2}\) 必须且只需 \(u_{1}\) 与 \(u_{2}\) 在由 \(\mathbf{C}^{n}\) 上以 A 为系数的多项式函数在 K 上的限制所组成的集合上相一致。

推论 2. --- 设 E 是一个 Banach 空间。设 K 是 \(\mathbf{C}^{n}\) 的一个多项式凸紧子集。设 P 是由 \(\mathbf{C}^{n}\) 上取值于 E 的多项式函数在 K 的某个邻域中的芽所组成的集合。则 \(\mathcal{O}(\mathrm{K};\mathrm{E})\) 的每个元素都是 P 中元素所成序列的极限。

给 E 赋以由 ab = 0（对 E 中所有的 a 与 b）定义的乘法 (I, p.~17, 例 1)。这是一个交换 Banach 代数。设 A 是由 E 添加单位元所得到的交换幺 Banach 代数。由于典范映射 \(\mathcal{O}(\mathrm{K};\mathrm{A}) \to \mathcal{O}(\mathrm{K};\mathrm{E})\) 连续，该断言由应用于代数 A 的 Oka-Weil 定理得出。

对 n = 1 与 A = C，我们还有下述结果，它将由 I, p.~150, 推论 2 加以精确化。

定理 3 (Runge). --- 设 K 是 C 的一个紧子集，Q 是在 K 的某个邻域中全纯的有理函数的芽所组成的集合。则 Q 在 \(\mathcal{O}(\mathrm{K})\) 中稠密。

按 \(\mathcal{O}(\mathrm{K})\) 上拓扑的定义，只需证明：对 K 的每个开邻域 U 以及 U 的每个紧子集 L，每个函数 \(f \in \mathcal{O}(\mathrm{U})\) 都是在 L 上连续的有理函数的极限。可以假设 L 是 K 的一个紧邻域。

设 Q' 是 C 上在 L 上连续的有理函数在 L 上的限制所组成的集合，并设 C 是 Q' 在 \(\mathcal{C}(\mathrm{L})\) 中的闭包。这是一个无根的代数。

设 \(z \in \mathrm{C}\) 是 L 的恒等映射。则 C 是 \(\mathcal{C}(\mathrm{L})\) 的由 \(z\) 生成的闭全子代数 (I, p.~6, 引理 2)。因此我们有 \(\mathrm{Sp}_{\mathrm{C}}(z) = \mathrm{Sp}_{\mathcal{C}(\mathrm{L})}(z) = \mathrm{L}\)。于是可以构造 C 中的元素 \(c = \Theta_z(f)\)。由于 C 无根，把 I, p.~66, 推论 1 应用于 C 的特征标 \(g \mapsto g(w)\)（对一切 \(w \in \mathrm{L}\)）就表明 \(c\) 与 \(f\) 在 L 上的限制重合。按 C 的定义，这就证明了 \(f|_{\mathrm{L}}\) 是 \(Q'\) 的元素在 L 上的一个一致极限，从而完成了定理的证明。

\subsection{全纯函数演算的存在性与唯一性}

假设 A 是一个复交换幺 Banach 代数。

定义 2. --- 设 \(n \geqslant 1\) 是一个整数，\(\boldsymbol{a} \in \mathrm{A}^{n}\)。设 U 是 \(\mathrm{Sp}^{n}(\boldsymbol{a})\) 的一个开邻域。我们称族 \(\boldsymbol{a}'\) 是 \((\boldsymbol{a}, \mathrm{U})\) 的一个包络，如果 \(\boldsymbol{a}' \in \mathbf{C}^{n+p}\) 延拓 \(\boldsymbol{a}\)，并且 \(\mathrm{U} \times \mathbf{C}^{p}\) 包含 \(\mathrm{Sp}^{n+p}(\boldsymbol{a}')\) 的多项式凸包络。

引理 11. --- 设 \(n \geqslant 1\) 是一个整数。设 \(\boldsymbol{a} \in \mathrm{A}^{n}\)。对 \(\mathrm{Sp}^{n}(\boldsymbol{a})\) 的每个开邻域 U，存在 \((\boldsymbol{a}, \mathrm{U})\) 的一个包络。

设 \((a_{\lambda})_{\lambda\in\Lambda}\) 是 A 的元素所成的一个族，它延拓族 a 并且拓扑生成幺 Banach 代数 A。设 \(\pi\) 是从 \(C^{\Lambda}\) 到 \(C^{n}\) 上的典范投影，并设 \(U'=\pi^{-1}(U)\)。则 \(U'\) 是 \(\mathrm{Sp}^{\Lambda}((a_{\lambda}))\) 的一个邻域，且 \(\mathrm{Sp}^{\Lambda}((a_{\lambda}))\) 是多项式凸的 (I, p.~44, 引理 4)。由 I, p.~47, 引理 6，存在一个有限子集 \(\Lambda_{0}\)（含于 \(\Lambda\) 且包含 \(\{1,2,\ldots,n\}\)），使得 \(\mathrm{pr}_{\Lambda_{0}}(\mathrm{U}')\) 包含 \(\mathrm{pr}_{\Lambda_{0}}(\mathrm{Sp}^{\Lambda}((a_{\lambda})_{\lambda\in\Lambda}))=\mathrm{Sp}^{\Lambda_{0}}((a_{\lambda})_{\lambda\in\Lambda_{0}})\) 的多项式凸包 S。设 \(p\geqslant0\) 是使 \(\Lambda_{0}\) 的基数为 \(n+p\) 的整数，并设 j 是从 \(\{1,\ldots,n+p\}\) 到 \(\Lambda_{0}\) 上的、在 \(\{1,\ldots,n\}\) 上与恒等映射重合的双射。由于 S 的投影包含于 U，族 \((a_{j(k)})_{1\leqslant k\leqslant n+p}\) 是 \((\boldsymbol{a},\mathrm{U})\) 的一个包络。

命题 5. --- 诸映射 \(\Theta_{\boldsymbol{a}}\)（对 \(n \geqslant 1\) 及 \(\boldsymbol{a} \in \mathrm{A}^{n}\)）的给定，构成 A 上的一个全纯函数演算，也就是说，I, p.~51 的条件 (CF1)、(CF2) 与 (CF3) 均得到满足。

设 \(n \geqslant 1\) 是一个整数，\(\boldsymbol{a} = (a_{1}, \ldots, a_{n}) \in \mathrm{A}^{n}\)。映射 \(\Theta_{\boldsymbol{a}}\) 满足 \(\Theta_{\boldsymbol{a}}(z_{i}) = a_{i}\)（对每个满足 \(1 \leqslant i \leqslant n\) 的 i；由 I, p.~64, 引理 10），这就证明了性质 (CF2)。I, p.~63, 引理 9 蕴涵诸映射 \(\Theta_{\boldsymbol{a}}\) 的性质 (CF3)。

映射 \(\Theta_{\pmb{a}}\) 是 A-线性的且连续的 (I, p.~61, 第 5 小节)。它满足 \(\Theta_{\pmb{a}}(1) = 1\) (I, p.~63, 引理 9)。为验证条件 (CF1)，只需再证明 \(\Theta_{\pmb{a}}\) 是一个代数态射。为此，我们将证明：\(\Theta_{\pmb{a}}^{\mathrm{U}}\)（对 \(\mathrm{Sp}^{n}(\pmb{a})\) 的每个开邻域 U）都是代数态射。

先设 U 包含 \(\mathrm{Sp}^{n}(\boldsymbol{a})\) 的多项式凸包 K。设 \(f_{1}\) 与 \(f_{2}\) 是 \(\mathcal{O}(\mathrm{U};\mathrm{A})\) 的元素。存在多项式函数的序列 \((f_{1,k})\)（相应地，\((f_{2,k})\)），它收敛于 \(f_{1}\)（相应地，收敛于 \(f_{2}\)）——收敛是在 \(\mathcal{O}(\mathrm{K};\mathrm{A})\) 中进行的 (I, p.~68, 定理 2)，因而在 \(\mathcal{O}(\mathrm{Sp}^{n}(\boldsymbol{a});\mathrm{A})\) 中收敛。对每个整数 \(k\)，有

\[
\Theta_ {\boldsymbol {a}} ^ {\mathrm{U}} (f _ {1, k}) \Theta_ {\boldsymbol {a}} ^ {\mathrm{U}} (f _ {2, k}) = \Theta_ {\boldsymbol {a}} ^ {\mathrm{U}} (f _ {1, k} f _ {2, k})
\]

（由 I, p.~64, 引理 10），取极限便得 \(\Theta_{a}^{\mathrm{U}}(f_{1})\Theta_{a}^{\mathrm{U}}(f_{2}) = \Theta_{a}^{\mathrm{U}}(f_{1}f_{2})\)。

考察一般情形。设 \(a' \in C^{n+p}\) 是 \((\boldsymbol{a}, U)\) 的一个包络（引理 11），\(\pi: U \times C^{p} \to U\) 是典范投影。由于 \(U \times C^{p}\) 包含 \(\mathrm{Sp}^{n+p}(\boldsymbol{a}')\) 的多项式凸包络，有

\[
\Theta_ {\boldsymbol {a} ^ {\prime}} ^ {\mathrm{U} \times \mathbf {C} ^ {p}} (f _ {1} \circ \pi) \Theta_ {\boldsymbol {a} ^ {\prime}} ^ {\mathrm{U} \times \mathbf {C} ^ {p}} (f _ {2} \circ \pi) = \Theta_ {\boldsymbol {a} ^ {\prime}} ^ {\mathrm{U} \times \mathbf {C} ^ {p}} (f _ {1} f _ {2} \circ \pi)
\]

（对 \(f_1\) 与 \(f_2\)——均属于 \(\mathcal{O}(\mathrm{U};\mathrm{A})\)——由第一种情形得到）。又因对每个函数 \(f \in \mathcal{O}(\mathrm{U};\mathrm{A})\)，有 \(\Theta_{\boldsymbol{a}'}^{\mathrm{U}\times \mathbf{C}^p}(f\circ \pi) = \Theta_{\boldsymbol{a}}^{\mathrm{U}}(f)\)（前面已证的条件 (CF3)），结论成立，从而条件 (CF1) 成立。

现在我们可以证明 I, p.~51 的定理 1。命题 5 表明映射族 \((\Theta_{a})_{a}\) 是 A 上的一个全纯函数演算。因此只需再确立 A 上全纯函数演算的唯一性。

设 \((\Psi_{a})_{a}\) 是一个映射族，它对每个整数 \(n \geqslant 1\) 以及每个 \(a \in A^{n}\) 有定义，并满足 A 上全纯函数演算的条件 (CF1)、(CF2)、(CF3) (I, p.~51)。只需证明：对每个整数 \(n \geqslant 1\)、每个 \(a \in A^{n}\) 以及 a 的每个开邻域 U，我们有 \(\Theta_{a}^{U} = \Psi_{a}^{U}\)。

设 \(n \geqslant 1\)，\(\boldsymbol{a} = (a_{1}, \ldots, a_{n}) \in \mathrm{A}^{n}\)。设 U 是 \(\mathrm{Sp}^{n}(\boldsymbol{a})\) 的一个开邻域。先设 U 包含 \(\mathrm{Sp}^{n}(\boldsymbol{a})\) 的多项式凸包 K。由性质 (CF1) 与 (CF2)，态射 \(\Theta_{a}^{U}\) 与 \(\Psi_{a}^{U}\) 在多项式函数上相一致。由 I, p.~68, 定理 2 的推论以及连续性性质 (CF1)，这两个态射因此相等。

我们来证明一般情形。设 \(a' \in C^{n+p}\) 是 \((\boldsymbol{a}, U)\) 的一个包络，\(\pi: U \times C^{p} \to U\) 是典范投影。我们有

\[
\Theta_ {\boldsymbol {a}} ^ {\mathrm{U}} = \Theta_ {\boldsymbol {a} ^ {\prime}} ^ {\mathrm{U} \times \mathbf {C} ^ {p}} \circ \pi^ {*} = \Psi_ {\boldsymbol {a} ^ {\prime}} ^ {\mathrm{U} \times \mathbf {C} ^ {p}} \circ \pi^ {*} = \Psi_ {\boldsymbol {a}} ^ {\mathrm{U}}
\]

（由性质 (CF3) 及前面的情形）。这就完成了 I, p.~51 定理 1 的证明。

注意，I, p.~68 的定理 2 还蕴涵下述唯一性结果：

命题 6. --- 设 \(\boldsymbol{a} \in \mathrm{A}^n\)。我们假设 \(\mathrm{Sp}^n(\boldsymbol{a})\) 是多项式凸的。设 \(z_1, \ldots, z_n\) 是 \(\mathrm{Sp}^n(\boldsymbol{a})\) 的某个邻域中的、\(\mathbf{C}^n\) 上坐标函数的芽。则映射 \(\Theta_{\boldsymbol{a}}\) 是唯一的连续幺代数态射 \(\varphi\)：从 \(\mathcal{O}(\mathrm{Sp}^n(\boldsymbol{a}); \mathrm{A})\) 到 A 中，使得 \(\varphi(z_1) = a_1, \ldots, \varphi(z_n) = a_n\)。

I, p.~64 的引理 10 与 I, p.~65 的命题 2 的推论为全纯函数演算论证了下述记号的合理性。设 \(n \geqslant 1\)，\(\boldsymbol{a} \in \mathrm{A}^n\)。对每个芽 \(f \in \mathcal{O}(\mathrm{K};\mathrm{A})\)（相应地，对每个全纯函数 \(f \in \mathcal{O}(\mathrm{U};\mathrm{A})\)——它定义在 \(\mathrm{Sp}^n (\boldsymbol{a})\) 的某个开邻域 U 上），我们令

\[
f (\boldsymbol {a}) = \Theta_ {\boldsymbol {a}} (f).\tag{6}
\]

依性质 (CF1) 与 (CF2)，当 f 是多项式时，这一记号与 A, IV, p.~4, n°3 中引入的记号相容。于是 I, p.~51 的性质 (CF2) 与 (CF3) 可以写成

\[
z _ {i} (\boldsymbol {a}) = a _ {i}, \quad 1 \leqslant i \leqslant n, \quad (f \circ \pi_ {m, n}) (\boldsymbol {a}) = f (\pi_ {m, n} (\boldsymbol {a})).
\]

\subsection{函数演算中的代入}

采用上面引入的记号，I, p.~66 的推论 1 与 I, p.~67 的推论 2 的结论分别成为

\[
\chi (g (\boldsymbol {a})) = g \left(\chi \left(a _ {1}\right), \dots , \chi \left(a _ {n}\right)\right), \quad \operatorname{Sp} (g (\boldsymbol {a})) = g \left(\operatorname{Sp} ^ {n} (\boldsymbol {a})\right)
\]

其中 \(f \in \mathcal{O}(\mathrm{Sp}^n(\boldsymbol{a});\mathrm{A})\)，\(\chi \in \mathrm{X}(\mathrm{A})\)，\(g \in \mathcal{O}(\mathrm{Sp}^n(\boldsymbol{a}))\)。

现在我们来证明一个更一般的代入性质。

定理 4. --- 设 A 是一个交换幺复 Banach 代数，设 \(n \geqslant 1\) 是一个整数，且 \(\boldsymbol{a} = (a_1, \ldots, a_n) \in \mathrm{A}^n\)。设 \(\boldsymbol{f} = (f_1, \ldots, f_p)\)，其中 \(f_1, \ldots, f_p\) 是 \(\mathcal{O}(\mathrm{Sp}^n(\boldsymbol{a}))\) 的元素。\(\mathrm{Sp}^n(\boldsymbol{a})\) 在映射 \(\boldsymbol{z} \mapsto \boldsymbol{f}(\boldsymbol{z}) = (f_1(\boldsymbol{z}), \ldots, f_p(\boldsymbol{z}))\) 下的像等于 \(\mathrm{Sp}^p(\boldsymbol{f}(\boldsymbol{a}))\)。

对每个 \(g \in \mathcal{O}(\mathrm{Sp}^{p}(\boldsymbol{f}(\boldsymbol{a}));\mathrm{A})\)，复合芽 \(g \circ f\) 是 \(\mathcal{O}(\mathrm{Sp}^{n}(\boldsymbol{a});\mathrm{A})\) 的元素，并且有 \(g(\boldsymbol{f}(\boldsymbol{a})) = (g \circ \boldsymbol{f})(\boldsymbol{a})\)。

关于 \(\mathrm{Sp}^{n}(\boldsymbol{a})\) 的像的第一个断言由 I, p.~67 的推论 2 得出。为证明第二个断言，我们将使用下述引理。

引理 12. --- 设 K 是 \(\mathrm{Sp}^p (\pmb {f}(\pmb {a}))\) 的多项式凸包。我们有 \(g(\pmb {f}(\pmb {a})) = (g\circ \pmb {f})(\pmb {a})\)（对每个芽 \(g\in \mathcal{O}(\mathrm{K};\mathrm{A})\)）。

设 \(\Psi\) 是从 \(\mathcal{O}(\mathrm{K};\mathrm{A})\) 到 A 的映射，满足 \(\Psi(g) = (g\circ f)(\boldsymbol{a})\)。它是一个连续的保单位元态射，满足 \(\Psi(z_j) = f_j(\boldsymbol{a})\)，其中 \(z_{j}\) 是第 \(j\) 个坐标函数（在 \(\mathbf{C}^p\) 上）的芽。因此，当 \(g\) 是多项式函数的芽时，有 \(\Psi(g) = g(\boldsymbol{f}(\boldsymbol{a}))\)。由 I, p.~68 的定理 2，此公式对所有 \(g\in \mathcal{O}(\mathrm{K};\mathrm{A})\) 仍然成立。

现在来证明该定理。设 V 是 \(\mathrm{Sp}^{p}(\boldsymbol{f}(\boldsymbol{a}))\) 的开邻域，\(\widetilde{g}\in\mathcal{O}(V;A)\) 是一个全纯函数，它在 \(\mathrm{Sp}^{p}(\boldsymbol{f}(\boldsymbol{a}))\) 的某邻域中的芽等于 g。设 \(b\in C^{p+q}\) 是 \((\boldsymbol{f}(\boldsymbol{a}),\mathrm{V})\) 的一个包络（I, p.~70 的引理 11），\(\pi\colon V\times C^{q}\to V\) 是典范投影。

设 \(\widetilde{f}_1, \ldots, \widetilde{f}_p\) 是以 \(f_1, \ldots, f_p\) 为芽的全纯函数，并设 U 是 \(\mathrm{Sp}^n(\boldsymbol{a})\) 的开邻域，使得 \((\widetilde{f}_1, \ldots, \widetilde{f}_p)(\mathrm{U}) \subset \mathrm{V}\)。设 \(\pi'\) 是从 \(\mathrm{U} \times \mathbf{C}^q\) 到 U 上的典范投影。用 \(z_{n+1}, \ldots, z_{n+q}\) 表示最后 \(q\) 个坐标函数（在 \(\mathbf{C}^{n+q}\) 上）。记 \(h = \widetilde{g} \circ (\widetilde{f}_1, \ldots, \widetilde{f}_p)\)，并令

\[
\boldsymbol {c} = (a _ {1}, \dots , a _ {n}, b _ {p + 1}, \dots , b _ {p + q}) \in \mathrm{A} ^ {n + q}.
\]

映射 \(\widetilde{g} \circ \pi\) 在开邻域 \(V \times C^q\)（即 \(Sp^{p+q}(b)\) 的多项式凸包 L 的开邻域）中全纯。将引理 12 应用于 \(c\) 以及下列函数在 L 的邻域中的芽

\[
(\widetilde {f} _ {1} \circ \pi^ {\prime}, \ldots , \widetilde {f} _ {p} \circ \pi^ {\prime}, z _ {n + 1}, \ldots , z _ {n + q}),
\]

以及 \(\widetilde{g}\circ\pi\) 的芽，便得

\[
(g \circ \pi) \left((f _ {1} \circ \pi^ {\prime}) (\boldsymbol {c}), \dots , (f _ {p} \circ \pi^ {\prime}) (\boldsymbol {c}), z _ {n + 1} (\boldsymbol {c}), \dots , z _ {n + q} (\boldsymbol {c})\right) = (h \circ \pi^ {\prime}) (\boldsymbol {c}).
\]

由于 \(\pi'(\boldsymbol{c}) = \boldsymbol{a}\)，有 \((h \circ \pi')(\boldsymbol{c}) = h(\boldsymbol{a})\) 与 \((f_i \circ \pi')(\boldsymbol{c}) = f_i(\boldsymbol{a})\)（对 \(1 \leqslant i \leqslant p\)；全纯函数演算的性质 (CF3)）。此外，有 \(z_{n+j}(\boldsymbol{c}) = b_{p+j}\)（对 \(1 \leqslant j \leqslant q\)）。由性质 (CF2)，有

\[
(g \circ \pi) (f _ {1} (\boldsymbol {a}), \dots , f _ {p} (\boldsymbol {a}), b _ {p + 1}, \dots , b _ {p + q}) = h (\boldsymbol {a}),
\]

再应用一次性质 (CF3)，便推得 \(g(f_{1}(\boldsymbol{a}), \ldots, f_{p}(\boldsymbol{a})) = h(\boldsymbol{a})\)。

\subsection{单变量的全纯函数演算}

定理 5. --- 设 A 是一个幺 Banach 代数，不必交换。设 a 是 A 的元素，z 是 C 的恒等函数在 SpA(a) 的某邻域中的芽。存在唯一的连续保单位元态射 \(\varphi_{a}\)：从 \(\mathcal{O}(\mathrm{Sp}_{\mathrm{A}}(a))\) 到 A 中，使得 \(\varphi_{a}(z)=a\)。

\(\varphi_{a}\) 的像包含在由 \(a\) 生成的 A 的闭全子代数中。特别地，它包含在 \(a\) 的双交换子中。

我们来证明态射 \(\varphi_{a}\) 的存在性。设 B 是由 \(a\) 生成的 A 的闭全子代数。它是交换的，并且 \(\mathrm{Sp}_{\mathrm{B}}(a) = \mathrm{Sp}_{\mathrm{A}}(a)\)（I, p.~5, 第 5 小节）。B 上的全纯函数演算的映射 \(\Theta_{a}\) 是从 \(\mathcal{O}(\mathrm{Sp}_{\mathrm{B}}(a))\) 到 B 的连续保单位元态射，满足 \(\Theta_{a}(z) = a\)（I, p.~51 的定理 1）。\(\Theta_{a}\) 与 B 到 A 中的典范内射的复合态射是一个连续保单位元态射 \(\varphi_{a}\)：从 \(\mathcal{O}(\mathrm{Sp}_{\mathrm{A}}(a))\) 到 A 中，它把 \(z\) 映为 \(a\)。

我们来证明唯一性。设 \(\varphi_{a}^{\prime}\) 是从 \(\mathcal{O}(\mathrm{Sp}_{\mathrm{A}}(a))\) 到 A 的连续保单位元态射，满足 \(\varphi_{a}^{\prime}(z)=a\)。那么 \(\varphi_{a}\) 与 \(\varphi_{a}^{\prime}\) 在 \(\mathrm{Sp}_{\mathrm{A}}(a)\) 的某邻域中的多项式的芽的集合上相一致，因而在 \(\mathrm{Sp}_{\mathrm{A}}(a)\) 的某邻域中的全纯有理分式的芽的集合上相一致。而这些芽在 \(\mathcal{O}(\mathrm{Sp}_{\mathrm{A}}(a))\) 中稠密（I, p.~69, 定理 3）。由此推出 \(\varphi_{a}=\varphi_{a}^{\prime}\)。

\(\varphi_{a}\) 的构造表明其像包含在交换子代数 B 中，而 B 包含在 a 的双交换子中（I, p.~6）。

如果代数 A 的根为零，态射 \(\varphi_{a}\) 的唯一性无需要求连续性即成立（参见 I, p.~40 的命题 9）。一般情形并非如此，参见 G. R. ALLAN, Embedding the algebra of formal power series in a Banach algebra, Proc. London Math. Soc. (3) 25 (1972), 329--340。

对每个 Banach 代数 A、A 的每个元素 \(a\) 以及每个芽 \(f \in \mathcal{O}(\mathrm{Sp}_{\mathrm{A}}(a))\)，我们用 \(f(a)\) 表示定理 5 中的元素 \(\varphi_{a}(f)\)。若 A 是交换 Banach 代数，则元素 \(f(a)\) 与交换 Banach 代数上的全纯函数演算所给出的元素 \(f(a)\) 相一致（I, p.~51 的定理 1）。

设 B 是由 \(a\) 生成的闭全子代数，于是 \(\mathrm{Sp}_{\mathrm{A}}(a) = \mathrm{Sp}_{\mathrm{B}}(a)\)。A 中的元素 \(f(a)\) 属于 B，并且与相对于代数 B 计算的元素 \(f(a)\) 相一致。

命题 7. --- 设 A 与 B 是幺 Banach 代数，\(\varphi\) 是从 A 到 B 的连续保单位元态射。设 \(a \in A\)。则 \(\mathrm{Sp}_{\mathrm{B}}(\varphi(a)) \subset \mathrm{Sp}_{\mathrm{A}}(a)\)，并且有 \(\varphi(f(a)) = f(\varphi(a))\)（对所有 \(f \in \mathcal{O}(\mathrm{Sp}_{\mathrm{A}}(a))\)）。特别地，对每个 \(\chi \in \mathsf{X}(\mathrm{A})\)，有 \(\chi(f(a)) = f(\chi(a))\)。

这由 I, p.~66 的命题 3 得出。

命题 8. --- 设 A 是一个幺 Banach 代数，\(a \in A\)。设 \(f \in \mathcal{O}(\mathrm{Sp}(a))\)。我们有 \(f(\mathrm{Sp}_{\mathrm{A}}(a)) = \mathrm{Sp}_{\mathrm{A}}(f(a))\)。此外，对每个 \(g \in \mathcal{O}(\mathrm{Sp}_{\mathrm{A}}(f(a)))\)，有 \(g \circ f \in \mathcal{O}(\mathrm{Sp}_{\mathrm{A}}(a))\) 且 \(g(f(a)) = (g \circ f)(a)\)。

这由定理 4 得出。

命题 9. --- 设 A 是一个幺 Banach 代数，\(a \in A\)。设 U 是 \(\mathrm{Sp}_{\mathrm{A}}(a)\) 的开邻域，\(f \in \mathcal{O}(\mathrm{U})\)。再设 V 是 \(\mathrm{Sp}_{\mathrm{A}}(a)\) 的包含在 U 中的紧邻域，使得 V 是 U 的一个具有定向边界 \(\partial\mathrm{V}\) 的块。

对每个整数 \(n \geqslant 0\)，映射 \(z \mapsto f(z)(z - a)^{-n-1}\) 在 \(\partial V\) 上连续；微分形式 \(z \mapsto f(z)(z - a)^{-n-1} dz\) 在 \(\partial V\) 上可积，并且有

\[
f ^ {(n)} (a) = \frac {n !}{2 i \pi} \int_ {\partial \mathrm{V}} f (z) (z - a) ^ {- n - 1} d z\tag{7}
\]

其中 \(f^{(n)}\in\mathcal{O}(\mathrm{U})\) 是 \(f\) 的第 n 阶导数。

对 \(n\) 用归纳法。当 \(n = 0\) 时，结论由 I, p.~61 的命题 1 得出。现设命题的断言对整数 \(n \geqslant 0\) 成立。设 \(g \in \mathcal{O}(\mathbf{C} - \mathrm{Sp}_{\mathrm{A}}(a); \mathrm{A})\) 是由 \(g(z) = (z - a)^{-n-1} f(z)\) 定义的全纯函数。微分形式 \(g'(z) dz = dg\) 是 \(\mathrm{C}^1\) 类的；由于块 V 是紧的，Stokes 公式（VAR, R2, p.~47, 11.2.3）蕴涵

\[
\int_ {\partial \mathrm{V}} g ^ {\prime} (z) d z = \int_ {\partial \mathrm{V}} d g = 0.
\]

由于 \(g'(z) = (z - a)^{-n-1} f'(z) - (n+1)(z - a)^{-n-2} f(z)\)，由此推出

\[
\int_ {\partial \mathrm{V}} f ^ {\prime} (z) (z - a) ^ {- n - 1} d z = (n + 1) \int_ {\partial \mathrm{V}} f (z) (z - a) ^ {- n - 2} d z.
\]

把归纳假设应用于 \(f'\)，便得

\[
\frac {2 i \pi}{n !} f ^ {(n + 1)} (a) = (n + 1) \int_ {\partial \mathrm{V}} f (z) (z - a) ^ {- n - 2} d z,
\]

这就是命题关于整数 \(n+1\) 的断言。证明完毕。

命题 10. --- 设 A 是一个幺 Banach 代数，U 是 C 的开子集。

\begin{enumerate}
\def\labelenumi{\alph{enumi})}
\tightlist
\item
  集合 \(\Omega\)——由这样的 \(a \in \mathrm{A}\) 组成：\(\operatorname{Sp}_{\mathrm{A}}(a) \subset \mathrm{U}\)——在 A 中是开的；
\item
  设 \(f \in \mathcal{O}(\mathrm{U})\)。映射 \(a \mapsto f(a)\)（从 \(\Omega\) 到 A 中）是全纯的，特别地是连续的。
\end{enumerate}

设 \(a \in \Omega\)。存在 \(\mathrm{Sp}_{\mathrm{A}}(a)\) 的包含在 U 中的紧邻域 V，它是 U 的一个块（VAR, R2, p.~46 与 p.~47, 11.1.3, d)）。

由于 \(a\) 的预解式在无穷远处趋于 0（I, p.~24 的定理 1, \(c\)）），映射 \(z \mapsto \| (z - a)^{-1}\|\) 在 \(\mathbf{C} - \mathring{\mathrm{V}}\) 上有界。用 M 表示它的上确界。若 \(h \in \mathrm{A}\) 满足 \(\| h\| \leqslant (2\mathrm{M})^{-1}\) 且 \(z \in \mathbf{C} - \mathring{\mathrm{V}}\)，则有

\[
z - (a + h) = (1 - h (z - a) ^ {- 1}) (z - a)
\]

并且 \(\|h(z-a)^{-1}\|\leqslant\frac{1}{2}\)，因而 \(z-(a+h)\) 可逆，其逆满足

\[
(z - (a + h)) ^ {- 1} = (z - a) ^ {- 1} \sum_ {n = 0} ^ {\infty} (h (z - a) ^ {- 1}) ^ {n}\tag{8}
\]

其中 \(\| (h(z - a)^{-1})^n\| \leqslant 2^{-n}\)（I, p.~22 的命题 2）。于是 \(\mathrm{Sp}_{\mathrm{A}}(a + h)\) 包含在 V 中，从而包含在 U 中，这就证明了 \(\Omega\) 在 A 中是开的。

设 \(f \in \mathcal{O}(\mathrm{U})\)。用 \(m\) 表示 \(|f(z)|\) 对 \(z \in \partial \mathrm{V}\) 的上确界。设 \(a \in \mathrm{A}\)。对每个 \(h \in \mathrm{A}\)（满足 \(\| h \| \leqslant (2\mathrm{M})^{-1}\)），有

\[
f (a + h) = \frac {1}{2 i \pi} \int_ {\partial \mathrm{V}} f (z) (z - (a + h)) ^ {- 1} d z
\]

（命题 9）。级数 (8) 在 V 的边界上一致收敛，因而

\[
f (a + h) = \sum_ {n = 0} ^ {+ \infty} f _ {a, n} (h)
\]

其中从 A 到 A 的映射 \(f_{a,n}\) 由下式定义

\[
f _ {a, n} (h) = \frac {1}{2 i \pi} \int_ {\partial \mathrm{V}} f (z) (z - a) ^ {- 1} (h (z - a) ^ {- 1}) ^ {n} d z.
\]

对每个 \(n \in N\)，函数 \(f_{a,n}\) 是 n 次连续齐次多项式函数。此外，由此可得

\[
\left\| f _ {a, n} (h) \right\| \leqslant \frac {m M}{\pi} \left(\int_ {\partial V} \| d z \|\right) 2 ^ {- (n + 1)}
\]

（INT, VI, §2, 第 3 小节, 命题 5）。因此级数 \(\sum_{n} f_{a,n}(h)\) 对 \(\|h\| \leqslant (2M)^{-1}\) 绝对收敛。这就证明了把 \(a\) 映为 \(f(a)\) 的映射在 \(\Omega\) 上全纯（VAR, R1, p.~26, 3.2.1）。

命题 11. --- 设 A 是一个幺 Banach 代数，\(a \in A\)，U 是 \(\mathrm{Sp}_{\mathrm{A}}(a)\) 的开邻域。用 \(\delta\) 表示 \(\mathrm{Sp}_{\mathrm{A}}(a)\) 到 \(\mathbf{C} - \mathbf{U}\) 的距离。设 \(f \in \mathcal{O}(\mathbf{U})\)。

\begin{enumerate}
\def\labelenumi{\alph{enumi})}
\item
  对每个满足 0 \textless{} η \textless{} δ 的实数 η，存在实数 C ≥ 0，使得 \(\|f^{(n)}(a)\| \leqslant \mathrm{Cn}!\eta^{-n}\)（对每个整数 \(n \in N\)）；
\item
  若 \(b \in A\) 与 \(a\) 交换且 \(\varrho(b) < \delta\)，则 \(\mathrm{Sp}_{\mathrm{A}}(a + b) \subset \mathrm{U}\)，并且
\end{enumerate}

\[
f (a + b) = \sum_ {n = 0} ^ {\infty} \frac {f ^ {(n)} (a)}{n !} b ^ {n},
\]

其中级数绝对收敛。

设 \(\eta\) 是满足 \(0 < \eta < \delta\) 的实数。记 \(\varepsilon = \delta - \eta > 0\)。设 K 是 \(\mathrm{Sp}_{\mathrm{A}}(a)\) 的紧邻域，它由 C 中到 \(\mathrm{Sp}_{\mathrm{A}}(a)\) 的距离 \(\leqslant \varepsilon / 2\) 的点组成。由于 \(f\) 在每个以属于 K 的点为中心、半径为 \(\eta + \varepsilon / 2\) 的开圆盘中全纯，由 Cauchy 不等式（VAR, R1, p.~29, 3.3.4），存在实数 \(C \geqslant 0\)，使得

\[
\sup _ {z \in K} \frac {| f ^ {(n)} (z) |}{n !} \leqslant \frac {C}{\eta^ {n}}
\]

对每个整数 \(n \geqslant 0\) 成立。于是把 I, p.~61 的命题 1 应用于 \(f^{(n)}\) 以及包含在 K 中的一个块 V，便得断言 a)。

设 \(b\) 是 A 中与 \(a\) 交换的元素，满足 \(\varrho(b) < \delta\)。用由 \(a\) 与 \(b\) 生成的闭全子代数 B 代替 A，该子代数满足 \(\mathrm{Sp}_{\mathrm{A}}(a) = \mathrm{Sp}_{\mathrm{B}}(a)\) 与 \(\mathrm{Sp}_{\mathrm{A}}(a + b) = \mathrm{Sp}_{\mathrm{B}}(a + b)\)，为证明 \(b\))，就可化为 A 为交换的情形。

由于 \(\varrho(b) < \delta\)，可以选取 \(\eta\) 使得 \(\varrho(b) < \eta < \delta\)。设 \(V_1\) 是 \(\mathbf{C}\) 中到 \(\mathrm{Sp}_{\mathrm{A}}(a)\) 的距离 \(< \delta - \eta\) 的点组成的集合，\(V_2\) 是以 0 为中心、半径为 \(\eta\) 的开圆盘（在 \(\mathbf{C}\) 中）。设 \(g\) 是映射 \((z_1, z_2) \mapsto z_1 + z_2\)：从 \(V_1 \times V_2\) 到 U 中。那么 \(h = f \circ g\) 是映射 \((z_1, z_2) \mapsto f(z_1 + z_2)\)：从 \(V_1 \times V_2\) 到 \(\mathbf{C}\) 中。我们有 \(\mathrm{Sp}_{\mathrm{A}}^{2}(a, b) \subset V_1 \times V_2\)，因而 \(\mathrm{Sp}_{\mathrm{A}}(a + b) \subset U\)（参见 I, p. 67 的推论 2），并且 \(f(a + b) = h(a, b)\)（由 I, p.~72 的定理 4）。现在，在空间 \(\mathcal{O}(V_1 \times V_2)\) 中，有

\[
h (z _ {1}, z _ {2}) = \sum_ {n \geqslant 0} \frac {f ^ {(n)} (z _ {1})}{n !} z _ {2} ^ {n},
\]

因此级数

\[
\sum_ {n \geqslant 0} \frac {f ^ {(n)} (a)}{n !} b ^ {n}
\]

在 A 中收敛，其和为 \(h(a,b)=f(a+b)\)。此外，由断言 (a)，这个级数绝对收敛。

\subsection{指数与对数}

我们用 exp 表示从 C 到 C 的复指数函数（FVR, III, p.~8, 定义 2）。它可微且满足 \(\exp' = \exp\)（FVR, III, p.~9, (26)）；因此它在 C 中全纯。设 A 是一个幺 Banach 代数，\(a\) 是 A 的元素。由 I, p.~65 的命题 2 与 FVR, III, p.~16 的公式 (9)，有

\[
\exp (a) = \sum_ {n = 0} ^ {\infty} \frac {a ^ {n}}{n !}.\tag{9}
\]

由于 \(\|a^{n}\|\leqslant\|a\|^{n}\)，我们看出 \(\|\exp(a)\|\leqslant\exp(\|a\|)\)，并且级数 (9) 在 A 的每个球中一致收敛。映射 \(a\mapsto\exp(a)\)（从 A 到 A）是全纯的（I, p.~76 的命题 10）。我们也时常用 \(e^{a}\) 表示 \(a\in A\) 的指数。

当 \(a\) 是 Banach 空间 E 的自同态时，像这样定义的指数 \(\exp(a)\)——在 Banach 代数 \(\mathcal{L}(\mathrm{E})\) 中——与 FVR, IV, p.~27, 定义 1 中定义的那个相一致，依据同前命题 7 (3)。

对 A 的每个与 a 可交换的元素 b，我们还有

\[
\exp (a + b) = \sum_ {n = 0} ^ {\infty} \frac {b ^ {n}}{n !} \exp (a),
\]

（I, p.~77 的命题 11），由此得

\[
\exp (a + b) = \exp (a) \cdot \exp (b).\tag{10}
\]

特别地，\(\exp(a)\) 可逆，且

\[
\exp (a) ^ {- 1} = \exp (- a).\tag{11}
\]

设 B 是由这样的 \(z \in C\) 组成的集合：\(-\pi < Jz < \pi\)。设 F 是 C 中区间 \(R_{-}\) 的余集。指数到 B 上的限制通过过渡到子空间诱导出从 B 到 F 上的一个双射（FVR, III, p.~10, n°7），其逆双射将记作 log。

若 \(a \in A\) 满足 \(\mathrm{Sp}_{\mathrm{A}}(a) \subset \mathrm{F}\)，就可以作出 A 的元素 \(\log(a)\)。我们有 \(\mathrm{Sp}_{\mathrm{A}}(\log(a)) \subset \mathrm{B}\)，并且

\[
\exp (\log (a)) = a\tag{12}
\]

（由 I, p.~75 的命题 8）。反之，设 \(b\) 是 A 中满足 \(\mathrm{Sp}_{\mathrm{A}}(b) \subset \mathrm{B}\) 的元素。我们有 \(\mathrm{Sp}_{\mathrm{A}}(\exp(b)) \subset \mathrm{F}\)，并且

\[
\log (\exp (b)) = b\tag{13}
\]

（同前）。

特别地，若 \(a \in A\) 满足 \(\varrho(a) < 1\)，则 \(\mathrm{Sp}_{\mathrm{A}}(1 - a) \subset \mathrm{F}\)，我们可以作出 \(\log(1 - a)\)。对 \(n \geqslant 1\)，第 \(n\) 阶导数（映射 \(z \mapsto \log(1 - z)\) 的）是 \(z \mapsto -(n - 1)! (1 - z)^{-n}\)。因此 \(z \mapsto \log(1 - z)\) 在点 0 处的幂级数展开为

\[
\log (1 - z) = - \sum_ {n = 1} ^ {\infty} \frac {z ^ {n}}{n},
\]

对 \(|z| < 1\) 成立（VAR, R1, p.~30, 3.3.9）。由 I, p.~65 的命题 2，我们得到

\[
\log (1 - a) = - \sum_ {n = 1} ^ {\infty} \frac {a ^ {n}}{n}.\tag{14}
\]

命题 12. --- 设 A 是一个交换幺 Banach 代数。指数映射的像是 A 的可逆元群 G 中含单位元的连通分支。

公式 (10) 与 (11) 证明 \(\exp(A)\) 是 G 的一个子群。由前文（见公式 (12)），这个子群包含以 1 为中心、半径为 1 的开球。因此它是 G 的开子群，从而是闭子群。此外，A 是连通的，且映射 \(a \mapsto \exp(a)\) 连续，所以 \(\exp(A)\) 连通。于是 \(\exp(A)\) 是 G 的含单位元的连通分支。

\subsection{特征标空间的分划}

命题 13. --- 设 A 是一个交换幺 Banach 代数。设 U₁ 与 U₂ 是 X(A) 的开子集，构成 X(A) 的一个分划。则 A 中存在唯一的幂等元 j，使得 Gelfand 变换 G(j) 在 U₁ 上等于 1，在 U₂ 上等于 0。

我们把空间 \(\mathsf{X}(\mathsf{A})\) 等同于 \(\mathbf{C}^{\mathrm{A}}\) 的一个紧子集，所借助的映射是 \(\chi \mapsto (\chi (a))_{a\in \mathrm{A}}\)（参见 I, p. 6, 第 6 小节以及 I, p. 29 定理 1 的推论）。子集 \(\mathrm{U}_1\) 与 \(\mathrm{U}_2\)（属于一致空间 \(\mathbf{C}^{\mathrm{A}}\)）是紧的且不相交。

由 TG, II, p.~31 的命题 4，存在 A 的有限子集 M 以及不相交的开集 \(\mathrm{V}_{1}\) 与 \(\mathrm{V}_{2}\)（属于 \(\mathbf{C}^{\mathrm{M}}\)），使得

\[
p (\mathrm{U} _ {1}) \subset \mathrm{V} _ {1}, \qquad p (\mathrm{U} _ {2}) \subset \mathrm{V} _ {2},
\]

其中 \(p\) 是 \(\mathbf{C}^{\mathrm{A}}\) 到 \(\mathbf{C}^{\mathrm{M}}\) 上的典范投影。

设 \(a_1, \ldots, a_n\) 是 M 的互不相同的元素，并把 \(\mathbf{C}^{\mathrm{M}}\) 等同于 \(\mathbf{C}^n\)。我们有 \(\mathrm{Sp}_{\mathrm{A}}^{n}(a_{1}, \ldots, a_{n}) \subset p(\mathsf{X}(\mathrm{A})) \subset \mathrm{V}_{1} \cup \mathrm{V}_{2}\)，因为 \(\mathrm{U}_{1} \cup \mathrm{U}_{2} = \mathsf{X}(\mathrm{A})\)。设 \(f\) 是 \(\mathrm{V}_{1} \cup \mathrm{V}_{2}\) 上的函数，在 \(\mathrm{V}_{1}\) 上等于 1，在 \(\mathrm{V}_{2}\) 上等于 0。我们有 \(f \in \mathcal{O}(\mathrm{V}_{1} \cup \mathrm{V}_{2})\)。令 \(j = f(a_{1}, \ldots, a_{n})\)。由于 \(f^{2} = f\)，有 \(j^{2} = j\)。由 I, p.~66 的推论 1，我们有 \(\chi(j) = 1\)（若 \(\chi \in \mathrm{U}_{1}\)），而 \(\chi(j) = 0\)（若 \(\chi \in \mathrm{U}_{2}\)），这就证明了所需幂等元的存在性。

另一方面，若 \(j_{1}\) 是 A 的幂等元，关系式 \(j^{2} = j\) 与 \(j_{1}^{2} = j_{1}\) 蕴涵 \((j - j_{1})(j + j_{1} - 1) = 0\)。若 \(\mathcal{G}(j_{1}) = \mathcal{G}(j)\)，则 \(j + j_{1} - 1\) 的 Gelfand 变换取值于 \(\{-1, 1\}\)，因而 \(j + j_{1} - 1\) 可逆（I, p.~37 的命题 6），于是 \(j = j_{1}\)。

推论. --- 设 A 是一个交换幺 Banach 代数。下列断言等价：

\begin{enumerate}
\def\labelenumi{\alph{enumi})}
\item
  特征标空间 X(A) 不是连通的；
\item
  A 存在异于 0 与 1 的幂等元。
\item
  代数 A 同构于两个非零 Banach 代数的乘积。
\end{enumerate}

该命题表明 \(a\)) 蕴涵 \(b\))。若 \(j\) 是 A 的幂等元，令 \(\mathrm{I}_1 = j\mathrm{A}\)，\(\mathrm{I}_2 = (1 - j)\mathrm{A}\)。那么 \(\mathrm{I}_1\) 与 \(\mathrm{I}_2\) 是 A 的闭理想，且 \(\mathrm{I}_1 + \mathrm{I}_2 = \mathrm{A}\)。若 \(j \notin \{0,1\}\)，则理想 \(\mathrm{I}_1\) 与 \(\mathrm{I}_2\) 都异于 A。另一方面，理想 \(\mathrm{I}_1\)（相应地，\(\mathrm{I}_2\)）是由 A 的元素 \(x\)（满足 \(jx = x\)，相应地，\((1 - j)x = x\)）组成的集合，因而 \(\mathrm{I}_1 \cap \mathrm{I}_2 = \{0\}\)。于是代数 A 等同于乘积 \(\mathrm{A} / \mathrm{I}_1 \times \mathrm{A} / \mathrm{I}_2\)。因此断言 \(b\)) 蕴涵 \(c\))。最后，若 A 同构于 \(\mathrm{A}_1 \times \mathrm{A}_2\)，则空间 \(\mathsf{X}(\mathsf{A})\) 等同于 \(\mathsf{X}(\mathsf{A}_1)\) 与 \(\mathsf{X}(\mathsf{A}_2)\) 的和空间（I, p.~6, 第 6 小节），因而 \(c\)) 蕴涵 \(a\))。

命题 14. --- 设 A 是一个无根的交换 Banach 代数。A 具有单位元，当且仅当 X(A) 是紧的。

条件是必要的（I, p.~29, 推论）。设 \(\mathsf{X}(\mathsf{A})\) 紧。设 \(\widetilde{\mathsf{A}}\) 是由 A 添加单位元而得到的 Banach 代数，并把 \(\mathsf{X}'(\mathsf{A})\) 等同于 \(\mathsf{X}(\widetilde{\mathsf{A}})\)。\(\mathsf{X}(\mathsf{A})\) 在 \(\mathsf{X}(\widetilde{\mathsf{A}})\) 中的余集化为一个特征标 \(\chi_0\)：它是 \(\widetilde{\mathsf{A}}\) 的特征标，其核为 A。集合 \(\mathsf{X}(\mathsf{A})\) 与 \(\{\chi_0\}\) 在 \(\mathsf{X}(\widetilde{\mathsf{A}})\) 中是开的。由命题 13，存在元素 \(j \in \mathsf{A}\)，使得 \(\chi(j) = 1\)（对 \(\chi \in \mathsf{X}(\mathsf{A})\)），且 \(\chi_0(j) = 0\)。于是我们有 \(j \in \mathsf{A}\)。

然后设 A 中的 \(x\)。有 \(\chi(jx) = \chi(x)\)（对所有 \(\chi \in \mathsf{X}(\mathsf{A})\)），由于 A 无根，故 \(jx = x\)。于是 \(j\) 是 A 的单位元。

命题 15. --- 设 A 是一个交换 Banach 代数，\(I_{1}\) 是 A 的理想，\(F_{1}\) 是由 \(\chi \in \mathsf{X}(A)\)（在 \(I_{1}\) 上取值为零）组成的集合。设 \(F_{2}\) 是 X(A) 中与 \(F_{1}\) 不相交的子集，对 Jacobson 拓扑是闭的，且对弱拓扑是紧的。则存在 \(u \in \text{I}_{1}\)，使得 \(\mathcal{G}(u) = 1\) 在 \(F_{2}\) 上成立。

设 \(I_{2}\) 是属于 \(F_{2}\) 的特征标的核的交。Banach 代数 A/ \(I_{2}\) 无根（I, p.~38 的命题 8）。由于 \(F_{2}\) 对 Jacobson 拓扑是闭的，X(A) 中在 \(I_{2}\) 上取值为零的元素只有 \(F_{2}\) 中的那些（参见 I, p.~13）。因此，赋以 X(A) 的弱拓扑所诱导的拓扑的 \(F_{2}\)，等同于赋以弱拓扑的 X(A/ \(I_{2}\) )（I, p.~9, 第 7 小节）。由于 \(F_{2}\) 弱紧，代数 A/ \(I_{2}\) 具有单位元（命题 14）。

于是我们有 \(I_{1} + I_{2} = A\)。事实上，若不然，\((I_{1} + I_{2})/I_{2}\) 将是真理想，从而包含在 \(A/I_{2}\) 的某个非零特征标的核中（I, p.~30, 定理 2）。这个特征标通过与典范投影 \(A \rightarrow A/I_{2}\) 复合，将定义 A 的一个非零特征标 \(\chi\)，它在 \(I_{1}\) 与 \(I_{2}\) 上取值为零，因而属于 \(F_{1} \cap F_{2}\)，与假设矛盾。

由于 \(I_{1} + I_{2} = A\)，存在 \(u \in I_{1}\)，其类在 \(A/I_{2}\) 中是 \(A/I_{2}\) 的单位元。那么 \(\chi(u) = 1\)（对所有 \(\chi \in F_{2}\)），证明完毕。

推论. --- 设 A 是一个交换 Banach 代数。设 \(F_{1}\) 与 \(F_{2}\) 是 \(\mathsf{X}(\mathsf{A})\) 的两个对 Jacobson 拓扑为闭的不相交子集。假设 \(F_{2}\) 弱紧。则存在 \(u \in A\)，使得 \(\mathcal{G}(u) = 1\) 在 \(F_{2}\) 上成立，且 \(\mathcal{G}(u) = 0\) 在 \(F_{1}\) 上成立。

\subsection{元素谱的分划}

设 A 是一个幺 Banach 代数，\(x \in A\)，且 \(K = \operatorname{Sp}_{\mathrm{A}}(x)\)。用 \(\Pi\) 表示 K 的在 K 中既开又闭的子集组成的集合。设 B 是由 x 生成的 A 的闭全子代数；它是交换的。

对每个 \(H \in \Pi\)，存在唯一的元素 \(f_{H}\)（属于 \(\mathcal{O}(K)\)），它在 H 的某邻域中等于 1，在 K---H 的某邻域中等于 0。我们令 \(j_{\mathrm{H}} = f_{\mathrm{H}}(x)\)。元素 \(j_{H}\) 是 A 的幂等元，称为与 x 和 H 相伴的幂等元，并且我们有下列公式：

\[
j _ {\mathrm{H} \cap \mathrm{H} ^ {\prime}} = j _ {\mathrm{H}} j _ {\mathrm{H} ^ {\prime}} = j _ {\mathrm{H} ^ {\prime}} j _ {\mathrm{H}}\tag{15}
\]

(H, H' ∈ Π)

\[
j _ {\mathrm {H\cup H^ {\prime}}} = j _ {\mathrm{H}} + j _ {\mathrm {H^ {\prime}}} - j _ {\mathrm {H^ {\prime}}} j _ {\mathrm{H}}\tag{16}
\]

\[
(\mathrm{H}, \mathrm{H} ^ {\prime} \in \Pi)
\]

\[
j _ {\emptyset} = 0, \qquad j _ {\mathrm{K}} = 1.
\]

设 \(H \in \Pi\)。我们定义 \(A_{H} = j_{H}Aj_{H}\)。这是 A 的一个闭子代数，以 \(j_{H}\) 为单位元（参见 I, p.~2 的引理 1）。我们还令 \(B_{H} = j_{H}Bj_{H}\) 以及 \(x_{H} = xj_{H} = j_{H}x = j_{H}xj_{H} \in B_{H}\)。

设 \(g_{H}\) 是 \(\mathcal{O}(K)\) 中由如下方式定义的元素：在 H 的邻域中 \(g_{\mathrm{H}}(z)=z\)，在 K---H 的邻域中 \(g_{\mathrm{H}}(z)=0\)。在 K 上有 \(g_{\mathrm{H}}(z)=f_{\mathrm{H}}(z)z\)，因此 \(x_{\mathrm{H}}=g_{\mathrm{H}}(x)\)。由此可得，若 \(H\neq K\)，则

\[
\mathrm{Sp} _ {\mathrm{A}} (x _ {\mathrm{H}}) = g _ {\mathrm{H}} (\mathrm{K}) = \mathrm{H} \cup \{0 \}.
\]

设 \(\lambda \in \mathbf{C} - \mathrm{H}\)。设 \(h_{\mathrm{H},\lambda}\) 是 \(\mathcal{O}(\mathrm{K})\) 中在 H 的某邻域中等于 \((\lambda - z)^{-1}\)、在 K-H 的某邻域中等于 0 的元素。我们有 \(h_{\mathrm{H},\lambda} = f_{\mathrm{H}}h_{\mathrm{H},\lambda}\) 以及 \((\lambda f_{\mathrm{H}} - g_{\mathrm{H}})h_{\mathrm{H},\lambda} = f_{\mathrm{H}}\)。若记 \(\mathrm{R_H}(x,\lambda) = h_{\mathrm{H},\lambda}(x)\)，则我们有 \(\mathrm{R_H}(x,\lambda) \in \mathrm{B_H}\)，并且

\[
\begin{array}{r} \mathrm {R_ {H}} (x, \lambda) (\lambda j _ {\mathrm{H}} - x _ {\mathrm{H}}) = (\lambda j _ {\mathrm{H}} - x _ {\mathrm{H}}) \mathrm {R_ {H}} (x, \lambda) = j _ {\mathrm{H}}, \\ \mathrm {R_ {H}} (x, \lambda) j _ {\mathrm {K-H}} = j _ {\mathrm {K-H}} \mathrm {R_ {H}} (x, \lambda) = 0. \end{array}\tag{17}
\]

特别地，\(\lambda \in \mathbf{C} - \mathrm{Sp}_{\mathrm{A_H}}(x_\mathrm{H})\)。

现在设 \(\lambda\in\mathrm{H}\)。假设 \(\lambda j_{\mathrm{H}}-x_{\mathrm{H}}\) 有逆 \(y\)（在 \(\mathrm{A}_{\mathrm{H}}\) 中）。利用公式 \(j_{\mathrm{H}}y=y\)（因为 \(y\in\mathrm{A}_{\mathrm{H}}\)）与 \(j_{\mathrm{K-H}}\mathrm{R}_{\mathrm{K-H}}(x,\lambda)=\mathrm{R}_{\mathrm{K-H}}(x,\lambda)\)（因为 \(\mathrm{R}_{\mathrm{K-H}}(x,\lambda)\in\mathrm{A}_{\mathrm{K-H}}\)），我们得到

\[
(\lambda - x) (y + \mathrm{R} _ {\mathrm{K-H}} (x, \lambda)) = (\lambda - x) (j _ {\mathrm{H}} y + j _ {\mathrm{K-H}} \mathrm{R} _ {\mathrm{K-H}} (x, \lambda)) =
\]

\[
(\lambda j _ {\mathrm{H}} - x j _ {\mathrm{H}}) y + (\lambda j _ {\mathrm {K-H}} - x j _ {\mathrm {K-H}}) \mathrm{R} _ {\mathrm {K-H}} (x, \lambda) = j _ {\mathrm{H}} + j _ {\mathrm {K-H}} = 1,
\]

（得益于应用于 K---H 的公式 (17)）。同样可以验证

\[
(y + \mathrm{R} _ {\mathrm{K-H}} (x, \lambda)) (\lambda - x) = 1.
\]

这证明 \(\lambda - x\) 在 A 中有逆 \(y + \mathrm{R}_{\mathrm{K - H}}(x, \lambda)\)，这是荒谬的。于是 \(\lambda \in \mathrm{Sp}_{\mathrm{A_H}}(x_{\mathrm{H}})\)。因此我们得出

\[
\mathrm{Sp} _ {\mathrm{A} _ {\mathrm{H}}} (x _ {\mathrm{H}}) = \mathrm{H}.\tag{18}
\]

特别地，若 H 非空，则幂等元 \(j_{H}\) 非零。

公式 (17) 与 (18) 证明：函数 \(\lambda \mapsto \mathrm{R}_{\mathrm{H}}(x,\lambda)\)（定义在 \(\mathbf{C} - \mathbf{H}\) 中）是 \(x_{\mathrm{H}}\) 相对于 \(\mathrm{A_H}\) 的预解式。

命题 16. --- 我们沿用前面的记号。设 \((\mathrm{H}_{i})_{1\leqslant i\leqslant n}\) 是 \(\mathrm{Sp}_{\mathrm{A}}(x)\) 的由 \(\Pi\) 中元素组成的一个分划。

\begin{enumerate}
\def\labelenumi{\alph{enumi})}
\item
  代数 B 典范地等同于代数 \(\mathrm{B_{H_1}}\times \dots \times \mathrm{B_{H_n}}\)
\item
  我们有 \(x_{\mathrm{H}_i}x_{\mathrm{H}_j} = 0\)（对 \(i\neq j\)），并且
\end{enumerate}

\[
x = x _ {\mathrm{H} _ {1}} + x _ {\mathrm{H} _ {2}} + \dots + x _ {\mathrm{H} _ {n}};
\]

\begin{enumerate}
\def\labelenumi{\alph{enumi})}
\setcounter{enumi}{2}
\tightlist
\item
  我们有
\end{enumerate}

\[
\mathrm{R} (x, \lambda) = \mathrm{R} _ {\mathrm{H} _ {1}} (x, \lambda) + \dots + \mathrm{R} _ {\mathrm{H} _ {n}} (x, \lambda)\tag{19}
\]

对所有 \(\lambda\in\mathbf{C}-\mathrm{Sp}_{\mathrm{A}}(x)\) 成立。特别地，若 \(\mathrm{H}\in\Pi\)，则预解式 \(\lambda\mapsto\mathrm{R}(x,\lambda)\) 在 \(\mathrm{H}\) 的邻域中等于 \(\mathrm{R}_{\mathrm{H}}(x,\lambda)\) 与一个全纯函数之和。

关系式 \(1 = j_{H_{1}} + \cdots + j_{H_{n}}\) 是 1 分解为 B 中两两正交的幂等元的分解，故代数 B 典范地等同于乘积代数 \(B_{H_{1}} \times \cdots \times B_{H_{n}}\)（A, I, p.~105, 命题 10）。

断言 \(b\)) 由函数 \(g_{\mathrm{H}_i}\) 的相应关系式得出；断言 \(c\)) 是 \(a\)) 与等式 \(\mathrm{R}(x_{\mathrm{H}},\lambda) = \mathrm{R}_{\mathrm{H}}(x,\lambda)\) 的推论。

命题 17. --- 设 \(\mu\) 是 \(\mathrm{Sp}_{\mathrm{A}}(x)\) 的一个孤立点。那么

\begin{enumerate}
\def\labelenumi{\alph{enumi})}
\tightlist
\item
  对每个 \(\lambda \in \mathbf{C} - \mathrm{Sp}_{\mathrm{A}}(x)\)，我们有
\end{enumerate}

\[
\mathrm{R} (x, \lambda) = \mathrm{R} _ {\{\mu \}} (x, \lambda) + \mathrm{R} _ {\mathrm{Sp} _ {\mathrm{A}} (x) - \{\mu \}} (x, \lambda);
\]

\begin{enumerate}
\def\labelenumi{\alph{enumi})}
\setcounter{enumi}{1}
\item
  函数 \(\lambda \mapsto \mathrm{R}_{\mathrm{Sp}_{\mathrm{A}}(x)-\{\mu\}}(x,\lambda)\) 在 \(\mathbf{C}-\mathrm{Sp}_{\mathrm{A}}(x)\) 中以及 \(\mu\) 的一个邻域中全纯；此外，函数 \(\lambda \mapsto \mathrm{R}_{\{\mu\}}(x,\lambda)\) 在 \(\mathbf{C}-\{\mu\}\) 中全纯；
\item
  我们有
\end{enumerate}

\[
\lim _ {n \to + \infty} \| (x - \mu) ^ {n} j _ {\{\mu \}} \| ^ {1 / n} = 0
\]

并且，对 \(\lambda \in \mathbf{C} - \{\mu\}\)，公式

\[
\mathrm{R} _ {\{\mu \}} (x, \lambda) = \sum_ {n = 0} ^ {\infty} (\lambda - \mu) ^ {- n - 1} (x - \mu) ^ {n} j _ {\{\mu \}}.\tag{20}
\]

前文已蕴涵断言 \(a\)) 与 \(b\))。我们来证明 \(c\))。通过把 \(x\) 换成 \(x-\mu\)，可化为 \(\mu=0\) 的情形。令 \(H = \{0\}\)；这是 \(\mathrm{Sp}_{\mathrm{A}}(x)\) 的既开又闭的子集。由公式 (18)，\(x_{H}\) 在 \(\mathrm{A}_{H}\) 中的谱是 \{0\}，因而 \(x_{H}\) 是拟幂零的，即 \(\|x^{n}j_{H}\|^{1/n}=\|(xj_{H})^{n}\|^{1/n}\) 当 \(n\) 趋于 + \(\infty\) 时趋于 0。此外，对 \(\lambda\neq0\)，在 A 中有

\[
\left(\lambda j _ {\mathrm{H}} - x _ {\mathrm{H}}\right) ^ {- 1} = \sum_ {n = 0} ^ {\infty} \lambda^ {- n - 1} x _ {\mathrm{H}} ^ {n}
\]

（I, p.~24 的定理 1, d)），由此得 (20)。

推论 1. --- 设 \(\mu\) 是 \(\mathrm{Sp}_{\mathrm{A}}(x)\) 的孤立点，\(p\) 是严格正整数。为使 \(\mu\) 成为阶为 \(p\) 的极点（对 \(x\) 的预解式而言；参见 VAR, R1, p.~30, 3.3.9），当且仅当 \((x - \mu)^{p-1} j_{\{\mu\}} \neq 0\) 且 \((x - \mu)^{p} j_{\{\mu\}} = 0\)。

推论 2. --- 设 \(\mu\) 是 \(\mathrm{Sp}_{\mathrm{A}}(x)\) 的孤立点。设 \(\Gamma\) 是一个开圆盘 \(\Delta\)（以 \(\mu\) 为中心）的定向边界，使得

\[
\operatorname{Sp} _ {\mathrm{A}} (x) \cap (\Gamma \cup \Delta) = \{\mu \}.
\]

则幂等元 \(j_{\{\mu \}}\)（与 \(x\) 和 \(\{\mu\}\) 相伴）由下式给出

\[
j _ {\{\mu \}} = \frac {1}{2 i \pi} \int_ {\Gamma} (z - x) ^ {- 1} d z.
\]

*换言之，幂等元 \(j_{\{\mu\}}\) 是在点 \(\mu\) 处 \(x.*\) 的预解式的留数。

对 \(z\in \mathbf{C} - \mathrm{Sp}_{\mathrm{A}}(x)\)，有

\[
(z - x) ^ {- 1} = \mathrm{R} (x, z) = \mathrm{R} _ {\{\mu \}} (x, z) + \mathrm{R} _ {\mathrm{H}} (x, z),
\]

其中 H = SpA(x) - \{μ\}（公式 (19)）。函数 z ↦ R\_H(x, z) 在 C-H 中以及 \{μ\} 的一个邻域中全纯（命题 17, b)），因而

\[
\frac {1}{2 i \pi} \int_ {\Gamma} \mathrm{R} _ {\mathrm{H}} (x, z) d z = 0
\]

（VAR, R2, p.~48, 11.2.5）。函数 \(z \mapsto \mathrm{R}_{\{\mu\}}(x,z)\) 是元素 \(j_{\{\mu\}}xj_{\{\mu\}}\)（属于幺代数 \(\mathrm{A}_{\{\mu\}}\)）的预解式。于是有

\[
j _ {\{\mu \}} = \frac {1}{2 i \pi} \int_ {\Gamma} \mathrm{R} _ {\{\mu \}} (x, z) d z
\]

（由应用于 \(A_{\{\mu\}}\) 以及 \(\Delta \cup \Gamma\) 的邻域中的常值函数 1 的 I, p.~75 的命题 9）。推论得证。

\subsection{完备可赋范的实或复代数中的全纯函数演算}

设 E 是一个实拓扑向量空间。拓扑向量空间 \(C \otimes E\)，即 E 的复化（EVT, II, p.~65），记作 \(\mathrm{E}_{(\mathbf{C})}\)，并且把 E 等同于 \(\mathrm{E}_{(\mathbf{C})}\) 的实拓扑向量子空间，所借助的映射是 \(x \mapsto 1 \otimes x\)。

命题 18. --- 复拓扑向量空间 \(E_{(\mathbf{C})}\) 是可赋范的（相应地，完备的），当且仅当 E 是可赋范的（相应地，完备的）。

\(\mathrm{E}_{(\mathbf{C})}\) 的底实拓扑向量空间同构于 \(E \times E\)。于是 \(\mathrm{E}_{(\mathbf{C})}\) 完备当且仅当 E 完备，且若 \(\mathrm{E}_{(\mathbf{C})}\) 可赋范，则 E 可赋范。

反之，设 E 可赋范。设 p 是定义 E 的拓扑的范数，B 是 p 的单位球。存在 \(\mathrm{E}_{(\mathbf{C})}\) 中 0 的闭均衡邻域 V，包含在 \(B+iB\) 中（EVT, II, p.~66）。于是集合 \(\lambda V\)（其中 \(\lambda\) 取遍 \(R_{+}^{*}\)）构成 \(\mathrm{E}_{(\mathbf{C})}\) 中 0 的基本邻域系。V 的规范函数是 \(\mathrm{E}_{(\mathbf{C})}\) 上的一个范数，它定义 \(\mathrm{E}_{(\mathbf{C})}\) 的拓扑，因而 \(\mathrm{E}_{(\mathbf{C})}\) 可赋范。

注记. --- 设 E 与 F 是 K 上的可赋范拓扑向量空间。从 E 到 F 中的连续线性映射的向量空间 \(\mathcal{L}(\mathrm{E};\mathrm{F})\)，赋以有界收敛拓扑，是一个可赋范拓扑向量空间（EVT, III, p.~14）。

设 E 与 F 是 R 上的可赋范拓扑向量空间。C-线性映射 \(\varphi\colon\mathcal{L}(\mathrm{E};\mathrm{F})_{(\mathbf{C})}\to\mathcal{L}(\mathrm{E}_{(\mathbf{C})};\mathrm{F}_{(\mathbf{C})})\)（由 \(\varphi(\lambda\otimes u)=\lambda u_{(\mathbf{C})}\) 定义）是复拓扑向量空间之间的同构。特别地，\(\mathrm{E}_{(\mathbf{C})}\) 的对偶等同于 E 的对偶的复化，赋范代数 \(\mathcal{L}(\mathrm{E}_{(\mathbf{C})})\) 等同于赋范代数 \(\mathcal{L}(\mathrm{E})\) 的复化。

设 S 是 C 中在复共轭下稳定的紧子集。考察 S 的邻域中取复值的全纯函数的芽所成的 C-代数 \(\mathcal{O}(S)\)，赋以 I, p.~49 第 1 小节中定义的复局部凸空间结构。若 U 是 S 在 C 中的开邻域，\(h: U \to C\) 是全纯函数，则 U 在复共轭下的像 V 是 S 在 C 中的开邻域，且 \(h^{*}: w \mapsto \overline{h(\overline{w})}\) 是 V 上的全纯函数。过渡到归纳极限，就得到一个连续对合 \(f\mapsto f^{*}\)（在代数 \(\mathcal{O}(\mathrm{S})\) 上）。特别地，我们有：

\[
(f + g) ^ {*} = f ^ {*} + g ^ {*} \quad (f g) ^ {*} = f ^ {*} g ^ {*} \quad (\lambda f) ^ {*} = \bar {\lambda} f ^ {*}
\]

对 \(f,g\)（属于 \(\mathcal{O}(\mathrm{S})\)）与 \(\lambda\)（属于 \(\mathbf{C}\)）成立。

我们用 \(\mathcal{O}_{\mathbf{R}}(\mathrm{S})\) 表示由芽 \(f \in \mathcal{O}(\mathrm{S})\)（满足 \(f = f^{*}\)）组成的集合。它是一个闭全子 \(\mathbf{R}\)-代数（\(\mathcal{O}(\mathrm{S})\) 中的）。

命题 19. --- 用 z 表示 \(\mathcal{O}(\mathrm{S})\) 中的、\(\mathbf{C}\) 的恒等映射的芽。则 \(\mathcal{O}_{\mathbf{R}}(\mathrm{S})\) 是最小的闭全子 \(\mathbf{R}\)-代数（在 \(\mathcal{O}(\mathrm{S})\) 中），且包含 \(z\)。

我们有 \(z^{*} = z\)，因而 z 属于 \(\mathcal{O}_{\mathbf{R}}(\mathrm{S})\)。设 B 是 \(\mathcal{O}(\mathrm{S})\) 的包含 z 的闭全子 R-代数。映射 \(f \mapsto f + f^{*}\)（从 \(\mathcal{O}(\mathrm{S})\) 到 \(\mathcal{O}_{\mathbf{R}}(\mathrm{S})\) 中）是连续且满射的，并且 S 的邻域中的全纯有理函数的芽的集合在 \(\mathcal{O}(\mathrm{S})\) 中稠密（I, p.~69 的定理 3）。因此，要证明 B 包含 \(\mathcal{O}_{\mathbf{R}}(\mathrm{S})\)，只需证明：若 f 是这样一个有理函数的芽，则有 \(f + f^{*} \in B\)。

存在 \(\mathbf{C}[X]\) 中的多项式 P 与 Q，使得 Q 在 S 的任何点上都不为零，并且有 \(f = \frac{\mathrm{P}(z)}{\mathrm{Q}(z)}\)。用 \(P^{*}\) 与 \(Q^{*}\) 表示把 P 与 Q 的系数换成其共轭而得到的多项式。于是有 \(\mathrm{P}(z)^{*} = \mathrm{P}^{*}(z)\) 与 \(\mathrm{Q}(z)^{*} = \mathrm{Q}^{*}(z)\)。由于 S 在复共轭下稳定，多项式 \(Q^{*}\) 在 S 的任何点都不为零。因此芽 \(\mathrm{Q}^{*}(z)\) 与 \((\mathrm{QQ}^{*})(z)\) 在 \(\mathcal{O}(\mathrm{S})\) 中可逆，并且

\[
f + f ^ {*} = \frac {\mathrm{P} (z)}{\mathrm{Q} (z)} + \frac {\mathrm{P} ^ {*} (z)}{\mathrm{Q} ^ {*} (z)} = \frac {(\mathrm{PQ} ^ {*} + \mathrm{P} ^ {*} \mathrm{Q}) (z)}{(\mathrm{QQ} ^ {*}) (z)}.
\]

由于多项式 PQ* + P*Q 与 QQ* 具有实系数，而 B 是 \(\mathcal{O}(\mathrm{S})\) 的包含 \(z\) 的全子 R-代数，故元素 \(f + f^{*}\) 属于 B。命题证明完毕。

设 A 是 R 上的一个完备可赋范幺代数。设 x 是 A 的元素。元素 \(1 \otimes x\)（属于代数 \(\mathrm{A}_{(\mathbf{C})}\)）的谱称为 x 的复谱，记作 \(\mathrm{Sp}_{\mathrm{A}_{(\mathbf{C})}}(x)\)。它与集合 R 的交正是 x 相对于 A 的谱 \(\mathrm{Sp}_{\mathrm{A}}(x)\)，后者有时称为 x 的实谱。复谱 \(\mathrm{Sp}_{\mathrm{A}_{(\mathbf{C})}}(x)\) 是 C 的紧子集，在复共轭下稳定。当代数 A 不化为 0 时它非空。

设 \(x\) 是 A 的元素。\(1 \otimes x \in \mathrm{A}_{(\mathbf{C})}\) 的谱半径等于谱半径 \(\varrho(x)\)（即 \(x\) 的谱半径）。它是最小的实数 \(r \geqslant 0\)，使得 \(|\lambda| \leqslant r\)（对所有 \(\lambda \in \mathrm{Sp}_{\mathrm{A}_{(\mathbf{C})}}(x)\)）。我们有

\[
\varrho (x) = \lim _ {n \to + \infty} \| x ^ {n} \| ^ {1 / n} = \inf _ {n > 0} \| x ^ {n} \| ^ {1 / n}
\]

对 A 上定义其拓扑的每个范数成立。事实上，可以假设 A 上的范数是 \(\mathrm{A}_{(\mathbf{C})}\) 上某个范数（它定义 \(\mathrm{A}_{(\mathbf{C})}\) 的拓扑）在 A 上的限制，并应用 I, p.~20 的命题 1。

我们用 \(u \mapsto \overline{u}\) 表示 R-代数 \(A_{(C)}\) 的把 \(\lambda \otimes a\) 映为 \(\overline{\lambda} \otimes a\) 的自同态。它是连续的。

引理. --- 对每个 \(f \in \mathcal{O}(\mathrm{Sp}_{\mathrm{A}_{(\mathbf{C})}}(x))\)，有 \(f^{*}(1 \otimes x) = \overline{f(1 \otimes x)}\)。

映射 \(f \mapsto f(1 \otimes x)\) 与 \(f \mapsto \overline{f^*(1 \otimes x)}\) 都是连续保单位元的 \(\mathbf{C}\)-代数同态：从 \(\mathcal{O}(\mathrm{Sp}(x))\) 到 \(\mathrm{A}_{(\mathbf{C})}\) 中，且都把 \(z\) 映为 \(1 \otimes x\)；因此它们相等（I, p.~74, 定理 5）。

命题 20. --- 对每个 \(f \in \mathcal{O}_{\mathbf{R}}(\mathrm{Sp}_{\mathrm{A}_{(\mathbf{C})}}(x))\)，存在 A 的唯一的元素 \(f(x)\)，使得 \(f(1 \otimes x) = 1 \otimes f(x)\)（在 \(\mathrm{A}_{(\mathbf{C})}\) 中）。映射 \(f \mapsto f(x)\)（从 \(\mathcal{O}_{\mathbf{R}}(\mathrm{Sp}_{\mathrm{A}_{(\mathbf{C})}}(x))\) 到 A 中）是唯一的连续保单位元 R-代数同态，它把 C 的恒等映射在 \(\mathcal{O}_{\mathbf{R}}(\mathrm{Sp}_{\mathrm{A}_{(\mathbf{C})}}(x))\) 中的芽映为 \(x\)。

我们记 \(S = \text{Sp}_{\text{A}(\mathbf{C})}(x)\)。由上面的引理，对每个芽 \(f \in \mathcal{O}_{\mathbf{R}}(\text{Sp}(x))\)，有 \(f(1 \otimes x) = \overline{f(1 \otimes x)}\)。由此得出第一个断言。用 \(z\) 表示 \(\mathcal{O}_{\mathbf{R}}(\text{S})\) 中的、\(\mathbf{C}\) 的恒等映射的芽。映射 \(f \mapsto f(x)\) 是从 \(\mathbf{R}\)-代数 \(\mathcal{O}_{\mathbf{R}}(\text{Sp}(x))\) 到 \(A\) 中的连续保单位元同态，它把 \(z\) 映为 \(x\)。由命题 19，这是唯一的这样一个同态，因为任何具有这些性质的态射在每个子 \(\mathbf{R}\)-代数（属于 \(\mathcal{O}(\text{S})\) 且包含 \(z\)）上都唯一确定。

设 \(f \in \mathcal{O}_{\mathbf{R}}(\mathrm{Sp}_{\mathrm{A}(\mathbf{C})}(x))\)。元素 \(f(x)\) 属于 A 的每个包含 \(x\) 的闭全子代数（命题 19），因而属于 \(x\) 在 A 中的双交换子。\(f(x)\) 的复谱等于 \(f(\mathrm{Sp}(x))\)（I, p.~75, 命题 8）。对所有 \(g \in \mathcal{O}_{\mathbf{R}}(f(\mathrm{Sp}_{\mathrm{A}(\mathbf{C})}(x)))\)，有 \(g \circ f \in \mathcal{O}_{\mathbf{R}}(\mathrm{Sp}_{\mathrm{A}(\mathbf{C})}(x))\) 且（同前）\(g \circ f(x) = g(f(x))\)。

设 U 是 C 的在复共轭下稳定的开子集。A 中复谱包含在 U 中的元素 x 组成的集合 \(\Omega\) 在 A 中是开的（I, p.~76, 命题 10）。设 f 是 U 上满足 \(f^{*} = f\) 的全纯函数。映射 \(x \mapsto f(x)\)（从 \(\Omega\) 到 A 中）是解析的（同前）。

设 A、B 是 R 上的完备可赋范幺结合代数，\(\varphi\colon A\to B\) 是连续的保单位元代数态射。设 \(x\in A\)。\(\varphi(x)\) 的复谱包含在 x 的复谱中，并且对每个 \(f\in\mathcal{O}_{\mathbf{R}}(\mathrm{Sp}_{\mathrm{A}(\mathbf{C})}(x))\)，有 \(f(\varphi(x))=\varphi(f(x))\)。这立即由复情形的相应命题得出（I, p.~75, 命题 8）。

\subsection{无单位元代数的情形}

设 A 是 K = R 或 C 上的完备可赋范代数，不必含单位元。用 ( \(\widetilde{A}, e\) ) 表示由 A 添加单位元而得到的幺代数。它是可赋范且完备的。

设 \(x\) 是 A 的元素。若 K = C，用 \(\mathrm{Sp}'(x) = \mathrm{Sp}_{\widetilde{\mathrm{A}}} (x)\) 表示 \(x\) 相对于 \(\widetilde{\mathrm{A}}\) 的谱，并考察芽 \(f \in \mathcal{O}(\mathrm{Sp}'(x))\)。若 K = R，用 \(\mathrm{Sp}'(x)\) 表示元素 \(x\)（属于 \(\widetilde{\mathrm{A}}\)）的复谱，并考察芽 \(f \in \mathcal{O}_{\mathbf{R}}(\mathrm{Sp}'(x))\)。在这两种情形，0 都属于 \(\mathrm{Sp}'(x)\)，并且元素 \(f(x)\)（属于 \(\widetilde{\mathrm{A}}\)）属于 A 当且仅当 \(f\) 满足 \(f(0) = 0\)。事实上，投影 \(\pi: \widetilde{\mathrm{A}} \to \mathrm{Ke}\) 是核为 A 的连续态射，并且有 \(\pi(f(x)) = f(\pi(x)) = f(0)\)。

\section{正则交换 Banach 代数}

在本节中，基域是 C。

\subsection{定义}

命题 1. --- 设 A 是一个交换 Banach 代数。下列条件等价：

\begin{enumerate}
\def\labelenumi{(\roman{enumi})}
\item
  \(\mathsf{X}(\mathsf{A})\) 上的弱拓扑与 Jacobson 拓扑相一致；
\item
  对每个 \(\chi \in \mathsf{X}(\mathrm{A})\) 与每个弱闭子集 \(\mathrm{F}\)（含于 \(\mathsf{X}(\mathrm{A})\)），使得 \(\chi \notin \mathrm{F}\)，存在 \(x \in \mathrm{A}\)，使得 \(\mathcal{G}(x)\) 在 \(\chi\) 处等于 1，在 \(\mathrm{F}\) 上等于 0；
\item
  对 X(A) 的每个弱紧子集 K 与每个弱闭子集 F，若 K ∩ F = ∅，则存在元素 \(x \in A\)，使得 \(\mathcal{G}(x)\) 在 K 上等于 1，在 F 上等于 0。
\end{enumerate}

设 \(M \subset \mathsf{X}(A)\)。说 M 对 Jacobson 拓扑是闭的，意思是：对每个 \(\chi \in \mathsf{X}(A) - \mathsf{M}\)，存在 \(x \in A\)，使得 \(\mathcal{G}(x)\) 在 M 上为零而不在 \(\chi\) 处为零（I, p.~39 的引理 2）。因此条件 (ii) 意味着 \(\mathsf{X}(A)\) 的每个弱闭子集对 Jacobson 拓扑是闭的，这表明 (ii) \(\Longrightarrow\) (i)。此外，(iii) \(\Longrightarrow\) (ii)，因为集合 \(\{\chi\}\) 在 \(\mathsf{X}(A)\) 中是弱紧的。最后，(i) \(\Longrightarrow\) (iii) 由 I, p.~81 的命题 15 的推论得出。

定义 1. --- 设 A 是一个交换 Banach 代数。若它满足命题 1 的等价条件，则称其为正则的。

注. --- 设 \(\tilde{\mathsf{A}}\) 是由 A 添加单位元 \(e\) 而得到的 Banach 代数。命题 1 的条件 (ii) 表明，若 \(\tilde{\mathsf{A}}\) 正则，则 A 正则。设 A 正则，我们来证明 \(\tilde{\mathsf{A}}\) 正则。考察 \(\mathsf{X}(\tilde{\mathsf{A}})\) 的不相交且弱闭（从而弱紧）的子集 F 与 F′，并构造 \(x \in \tilde{\mathsf{A}}\)，使得 \(\mathcal{G}(x)\) 在 F 上为零而在 F′ 上等于 1。设 \(\chi_0 \in \mathsf{X}(\tilde{\mathsf{A}})\) 是 A 上的零特征标。若 \(\chi_0 \notin \mathrm{F}'\)，则由命题 1 的条件 (iii) 以及关于 A 的假设，存在元素 \(x \in \mathsf{A}\)，使得 \(\mathcal{G}(x)\) 在 F 上为零而在 F' 上等于 1。若 \(\chi_0 \in \mathrm{F}'\)，则 \(\chi_0 \notin \mathrm{F}\)；因此存在元素 \(y \in \mathsf{A}\)，使得 \(\mathcal{G}(y)\) 在 F' 上为零而在 F 上等于 1。于是 \(x = e - y\)（属于 \(\tilde{\mathsf{A}}\)）具有所需的性质。

例. --- 我们回到 I, p.~17 第 2 小节的例子。

局部紧空间 X 上在无穷远处趋于 0 的连续复值函数所成的代数（I, p.~17 的例 3）是正则的（参见 I, p.~36, 例 1）。

 \(n\) 次可微且定义在 \([0, 1]\) 上的函数所成的代数（I, p.~18 的例 4）是正则的（参见 I, p.~36, 例 2）。

若 G 是交换局部紧群，\(\mu\) 是 G 上的一个 Haar 测度，则代数 \(L^{1}(G,\mu)\)（I, p.~19 的例 7）是正则的（参见 II, p.~219, 推论 2）。

在圆盘 \(|z| \leqslant 1\) 上连续且在其内部解析的函数所成的代数（I, p.~20 的例 9）不是正则的（参见 I, p.~193, 习题 6）。

命题 2. --- 设 A 是一个交换幺正则 Banach 代数。设 \(n \geqslant 1\) 是一个整数，\((\mathrm{U}_{1}, \ldots, \mathrm{U}_{n})\) 是 \(\mathsf{X}(\mathsf{A})\) 的一个开覆盖。存在 A 中和为 1 的元素 \(x_{1}, \ldots, x_{n}\)，使得 \(\operatorname{Supp}(\mathcal{G}(x_{i})) \subset \mathrm{U}_{i}\)（其中 \(i = 1, \ldots, n\)）。

我们对 n 用归纳法证明该命题。断言对 n = 1 成立。设 \(n \geqslant 2\)，并且断言对 n - 1 已被证明。

存在 \((\mathrm{V}_{1},\ldots,\mathrm{V}_{n})\)（它是 \(\mathsf{X}(\mathsf{A})\) 的开覆盖），使得对每个 i 有 \(\overline{V}_{i}\subset U_{i}\)。由归纳假设，存在元素 \(x,x_{3},\ldots,x_{n}\in A\)，使得 \(x+x_{3}+\cdots+x_{n}=1\)，且 \(\operatorname{Supp}(\mathcal{G}(x))\subset\mathrm{V}_{1}\cup\mathrm{V}_{2}\)，以及对 \(\operatorname{Supp}(\mathcal{G}(x_{i}))\subset\mathrm{V}_{i}\)，有 \(i\geqslant3\)。令 \(\mathrm{K}=\operatorname{Supp}(\mathcal{G}(x))\subset\mathrm{V}_{1}\cup\mathrm{V}_{2}\)。设 \(K_{1}\)（相应地，\(K_{2}\)）为 K 中不属于 \(V_{1}\)（相应地，\(V_{2}\)）的元素组成的集合。那么 \(K_{1}\) 与 \(K_{2}\) 是 K 的不相交紧子集。由于 Banach 代数 A 是正则的，故存在 \(y\in A\)，使得 \(\mathcal{G}(y)=1\) 在 \(K_{1}\) 上成立，且 \(\mathcal{G}(y)=0\) 在 \(K_{2}\) 上成立。于是 \(\mathcal{G}(xy)\) 在 \(\mathsf{X}(\mathsf{A})-\mathsf{K}\) 与 \(K_{2}\) 上为零，因而 \(\operatorname{Supp}\mathcal{G}(xy)\subset\overline{\mathrm{V}}_{2}\subset\mathrm{U}_{2}\)。同理，\(\mathcal{G}(x(1-y))\) 在 \(\mathsf{X}(\mathsf{A})-\mathsf{K}\) 与 \(K_{1}\) 上为零，因而 \(\operatorname{Supp}\mathcal{G}(x(1-y))\subset\overline{\mathrm{V}}_{1}\subset\mathrm{U}_{1}\)。于是元素 \(x_{1}=x(1-y)\)、\(x_{2}=xy\) 与 \(x_{3},\ldots,x_{n}\) 满足命题所述的性质。

推论 1. --- 设 A 是一个交换幺正则 Banach 代数，I 是 A 的理想，\(f: \mathsf{X}(\mathsf{A}) \to \mathbf{C}\) 是连续函数。假设对每个 \(\chi \in \mathsf{X}(\mathsf{A})\)，存在元素 \(y_{\chi} \in \mathrm{I}\)，使得在 \(f = \mathcal{G}(y_{\chi})\) 的一个邻域中有 \(\chi\)。则存在元素 \(y \in \mathrm{I}\)，使得 \(f = \mathcal{G}(y)\)。

由于 X(A) 紧，存在 X(A) 的有限开覆盖 \((\mathrm{U}_{1},\ldots,\mathrm{U}_{n})\)，以及 I 的元素 \(y_{1},\ldots,y_{n}\)，使得在 \(f=\mathcal{G}(y_{i})\) 上有 \(U_{i}\)。由命题 2，存在 A 中和为 1 的元素 \(x_{1},\ldots,x_{n}\)，使得对每个 i 有 \(\operatorname{Supp}(\mathcal{G}(x_{i}))\subset\mathrm{U}_{i}\)。令 \(y=x_{1}y_{1}+\cdots+x_{n}y_{n}\)。这是 I 中具有所需性质的元素。事实上，设 \(\chi\in\mathrm{X}(A)\)。对 \(1\leqslant i\leqslant n\)，我们有 \(\mathcal{G}(x_{i})(\chi)\mathcal{G}(y_{i})(\chi)=\mathcal{G}(x_{i})(\chi)f(\chi)\)，因为若 \(\mathcal{G}(y_{i})(\chi)=f(\chi)\)，则 \(\chi\in U_{i}\)，而若 \(\mathcal{G}(x_{i})(\chi)=0\)，则 \(\chi\notin U_{i}\)。由此得出

\[
\mathcal {G} (y) (\chi) = \sum_ {i = 1} ^ {n} \mathcal {G} (x _ {i}) (\chi) \mathcal {G} (y _ {i}) (\chi) = f (\chi) \sum_ {i = 1} ^ {n} \mathcal {G} (x _ {i}) (\chi) = f (\chi).
\]

推论 2. --- 设 A 是一个交换正则 Banach 代数，I 是 A 的理想，\(f\colon\mathsf{X}'(\mathrm{A})\to\mathbf{C}\) 是连续函数。假设对每个 \(\chi\in\mathsf{X}'(\mathrm{A})\)，存在元素 \(y_{\chi}\in\mathrm{I}\)，使得在 \(f=\mathcal{G}'(y_{\chi})\) 的一个邻域中有 \(\chi\)。则存在元素 \(y\in\mathrm{I}\)，使得 \(f=\mathcal{G}'(y)\)。

设 \(\tilde{\mathsf{A}}\) 是由 \(\mathsf{A}\) 添加单位元而得到的 Banach 代数。那么 \(\tilde{\mathsf{A}}\) 是正则的（注 1），且 \(\mathsf{X}'(\mathsf{A}) = \mathsf{X}(\widetilde{\mathsf{A}})\)；因此只需对 \(\tilde{\mathsf{A}}\) 与理想 I 应用推论 1。

若 I 是交换 Banach 代数的理想，回顾（参见 I, p.~30）：我们用 V(I) 表示核包含 I 的 \(\chi \in \mathsf{X}(\mathrm{A})\) 组成的集合，换言之，即这样的 \(\chi \in \mathsf{X}(\mathrm{A})\) 组成的集合：所有对应于 \(\mathcal{G}(x)\)（\(x \in \mathrm{I}\)）的函数都在其上为零。这是 \(\mathsf{X}(\mathrm{A})\) 的一个对 Jacobson 拓扑为闭的子集。

命题 3. --- 设 A 是一个交换正则 Banach 代数，I 是 A 的理想，K 是 X(A) 的与 V(I) 不相交的紧子集。存在元素 \(x \in I\)，使得 \(\mathcal{G}(x) = 1\) 对 K 中的每个 \(x\) 成立。

这是 I, p.~81 的命题 15 的一个特殊情形，只需注意到 Jacobson 拓扑与 X(A) 上的弱拓扑相一致。
